\documentclass[twoside,english]{paper}
\usepackage[T1]{fontenc}
\usepackage[latin9]{inputenc}
\usepackage{geometry}
\geometry{verbose}
\usepackage{babel}
\usepackage{mathrsfs}
\usepackage{amsmath}
\usepackage{amsthm}
\usepackage{amssymb}
\usepackage{booktabs}
\usepackage{enumitem}
\usepackage{longtable}
\usepackage{float}
\usepackage[pdfusetitle,
 bookmarks=true,bookmarksnumbered=false,bookmarksopen=false,
 breaklinks=false,pdfborder={0 0 1},backref=false,colorlinks=false]
 {hyperref}

\makeatletter
%%%%%%%%%%%%%%%%%%%%%%%%%%%%%% Textclass specific LaTeX commands.
% Save and disable \example as this might
% clash with theorems
\let\paperclassexample\example
\let\endpaperclassexample\endexample
\let\example\relax
% If theorems hasn't been loaded, restore example
\AtBeginDocument{%
  \@ifundefined{example}{\let\example\paperclassexample}{}
}
\numberwithin{equation}{section}
\numberwithin{figure}{section}
\newcommand{\lyxaddress}[1]{
	\par {\raggedright #1
	\vspace{1.4em}
	\noindent\par}
}

\makeatother

% ---------
% Theorem environments
% ---------
\newtheorem{theorem}{Theorem}[section]
\newtheorem{proposition}[theorem]{Proposition}
\newtheorem{lemma}[theorem]{Lemma}
\newtheorem{corollary}[theorem]{Corollary}
\theoremstyle{definition}
\newtheorem{definition}[theorem]{Definition}
\newtheorem{example}[theorem]{Example}
\theoremstyle{remark}
\newtheorem{remark}[theorem]{Remark}

% ---------
% Notation
% ---------
\newcommand{\OO}{\mathbb{O}}          % octonions
\newcommand{\RR}{\mathbb{R}}
\newcommand{\QQ}{\mathbb{Q}}
\newcommand{\ZZ}{\mathbb{Z}}
\newcommand{\FF}{\mathbb{F}}
\newcommand{\CC}{\mathbb{C}}
\newcommand{\HH}{\mathbb{H}}
\newcommand{\LL}{\Lambda}             % Leech lattice
\newcommand{\norm}[1]{\lVert #1 \rVert}
\newcommand{\ip}[2]{\langle #1, #2 \rangle}
\newcommand{\Nrm}{N}                  % norm form
\DeclareMathOperator{\GL}{GL}
\DeclareMathOperator{\Co}{Co}
\DeclareMathOperator{\Aut}{Aut}

% Verification scripts are cited in footnotes, hot-linked into the public
% research record.  Argument: the LaTeX-escaped file name.
\newcommand{\scriptref}[1]{\cite{repo}~%
  \href{https://github.com/koeplinger/leech\_alg/blob/main/python\_project/src/#1}%
  {\nolinkurl{python_project/src/#1}}}
\newcommand{\gapref}[1]{\cite{repo}~%
  \href{https://github.com/koeplinger/leech\_alg/blob/main/gap\_project/#1}%
  {\nolinkurl{gap_project/#1}}}

\setlength{\emergencystretch}{3em}

\begin{document}
\title{An order on the Leech lattice from a $\ZZ_3$-symmetric
triple-octonion product}
\subtitle{18 July 2026 (v6)}
\author{Jens K\"oplinger}
\maketitle

\lyxaddress{105 E Avondale Dr, Greensboro, NC 27403, USA \\
jenskoeplinger@gmail.com}

\begin{abstract}
We show that the Leech lattice~$\LL$ admits the structure of an order
in a 24-dimensional real algebra: the lattice is closed under a
bilinear product on~$\RR^{24}$.  The product is a
$\ZZ_3$-symmetric triple-octonion product, assembled
from nine octonion-multiplication blocks under cyclic $\ZZ_3$
permutation of three octonion factors.  Using Robert A.~Wilson's sublattice
description of~$\LL$ [J.~Algebra \textbf{322} (2009), 2186--2190],
closure turns on a transposition~$\sigma$ of
two imaginary basis units: $\sigma$ leaves the $E_8$ lattice~$L$
invariant and is an algebra isomorphism of the octonions, but it
moves Wilson's sublattices $L\bar s$ and~$Ls$.

Furthermore, using exact arithmetic (sampled statistics identified as
such) and computer proof, we offer
a compact exposition for the applied researcher: the algebraic properties of this
product, its idempotents and square-zero elements, its automorphisms, the
span of its image and the index of that span, and which permutation cycles,
beyond the transposition~$\sigma$, also close.  Reduced modulo~$2L$, the construction produces a Witt
decomposition of $L/2L$ into a complementary pair of Lagrangians, opening
the natural research program of which polarizations of $L/2L$ lift to a
closed product on~$\LL$.

We further provide a historical account from the original sources
(1923--1946) in a historical appendix.  Leonard E.~Dickson gave the
first verified maximal order of integral octonions (1923), with an
undercount of the total.  Johannes Kirmse (1924) derived a necessary form
met by thirty candidates, and hinted, without detail, at a further
narrowing; in our reconstruction it leaves eight, containing all seven
maximal orders, though the one he exhibited is not closed.  A transposition
provided by Richard H.~Bruck repaired that candidate (1946), and
H.\,S.\,M.~(``Donald'') Coxeter gave the first correct account of all seven
maximal orders.  The map~$\sigma$ of the present construction is exactly
this transposition, acting on a different sublattice.

This result was developed through a systematic human--AI collaboration;
the complete research record is publicly available.
\end{abstract}

\medskip
\noindent\textbf{Keywords:} Leech lattice; integral octonions; $E_8$
lattice; nonassociative algebra; non-unital order; computer-assisted
proof.

\smallskip
\noindent\textbf{MSC 2020:} Primary 11H56, 17A75; Secondary 11R52,
17D05, 68V05.

% ===============================================================
\section{Introduction}\label{sec:intro}
% ===============================================================

The Leech lattice~$\LL$ is the unique even unimodular lattice of
rank~24 with no vectors of squared norm~2
\cite{ConwaySloane1999}.  Its minimal shell $\operatorname{Min}(\LL)$
consists of 196{,}560 vectors of squared norm~8, and its automorphism
group is the Conway group~$\Co_0$, a double cover of the sporadic
simple group~$\Co_1$.

In rank~8, the $E_8$ lattice~$L$ carries an octonionic multiplication
under which it is closed.
Seven octonion-multiplication conventions on~$\RR^8$ have this
property; the resulting structures are in bijection with the seven
maximal orders of the integral octonions
\cite{Coxeter1946, ConwaySmith2003}.  Wilson
\cite[Section~3]{Wilson2009} showed that $\LL$ embeds naturally
in~$L^3$ via three sublattice conditions, raising the question: does
$\LL$ inherit an algebraic structure from the octonion product?

This question has been studied from several angles.  Geoffrey
M.~Dixon \cite{Dixon2010} constructed~$\LL$ using an octonionic
XY-product.  John C.~Baez and Greg Egan \cite{Baez2014} explored
Jordan-algebraic structures on~$\LL$ within the exceptional Jordan
algebra~$\mathfrak{h}_3(\OO)$, finding closure under a doubled
Jordan product.

The $\ZZ_3$-symmetric triple-octonion product on~$L^3$ uses the
same octonion multiplication in all three blocks with
$\ZZ_3$-symmetric cross-block routing.  It satisfies two of Wilson's
three sublattice conditions automatically.  The third, however,
generally fails: the sublattice it involves is not closed under
octonion products, including Wilson's.

The main result of this paper is that closure on~$\LL$ is restored
by twisting the octonion multiplication on each block by a
transposition~$\sigma$ of two imaginary basis units.  The map
$\sigma$ fixes the $E_8$ lattice~$L$ setwise and is an
octonion-algebra isomorphism, but it moves Wilson's sublattices
$L\bar s$ and~$Ls$ onto sublattices for which the required closure
inclusions do hold.

\begin{theorem}\label{thm:main}
Let\/ $\star$ denote the $\ZZ_3$-symmetric triple-octonion product
on~$\RR^{24} = \OO_1 \oplus \OO_2 \oplus \OO_3$
under $\ZZ_3$ cross-block routing (Definition~\ref{def:triple}).
Then $u \star v \in \LL$ for all $u, v \in \LL$ in Wilson's
representation~\cite{Wilson2009}.
\end{theorem}

\begin{corollary}\label{cor:order}
$(\LL, +, \star)$ is an order\footnote{The ambient algebra has no
multiplicative identity (Section~\ref{sec:properties}), so the term
``order'' is used here in the non-unital sense.} in the $\RR$-algebra
$(\RR^{24}, +, \star)$.
\end{corollary}

\begin{proof}
The product~$\star$ is bilinear (it is a sum of bilinear octonion
products).  By Theorem~\ref{thm:main}, $\star$ maps $\LL \times \LL$
into~$\LL$.
\end{proof}

\begin{remark}[Order-closure, not shell-closure]\label{rem:order-closure}
The closure statement $\LL \star \LL \subseteq \LL$ means that
products land somewhere in~$\LL$; it does not assert that the
minimal shell is preserved.  The product of two minimal vectors has
norm at least~16 (Remark~\ref{rem:product-norms}), never the minimal
value $\Nrm = 8$: such products always land on strictly higher
shells.  This is consistent with $(\LL, +, \star)$ being an
order in the ring-theoretic sense: a subring that is a free
$\ZZ$-module of rank equal to the ambient algebra's dimension.
That property is weaker than, and independent of, any shell-level
closure.
\end{remark}

The scope of this paper is deliberately bounded.  We establish
existence of the order, with a proof that is elementary (four
short lemmas, reducing to sublattice inclusions verified in exact
arithmetic).  Further structural readings are offered as starting points for further investigation in
Section~\ref{sec:properties} and Appendix~\ref{sec:appendix-tables}.
Nowhere do we claim or advertise these computer proofs as a substitute
for mathematical understanding: they are listed as machine-obtained
facts, to provide pointers for the applied reader.  Interest in earlier
versions of this work came from researchers in medical imaging,
cryptography, quantum computing, and phenomenology in the physics of
fundamental particles.

All results were developed through a systematic human--AI
collaboration, and the complete research record is public~\cite{repo}:
every prompt, computation, and verification, end to end
(Appendix~\ref{sec:methodology}).

% ===============================================================
\section{Preliminaries}\label{sec:prelim}
% ===============================================================

\subsection{The octonion algebra}

Let $\OO$ denote the real octonion algebra with basis
$\{e_0, e_1, \ldots, e_7\}$, where $e_0$ is the identity.
The multiplication is determined by the seven \emph{Fano triples}
\begin{equation}\label{eq:fano}
  (1,2,4),\; (2,3,5),\; (3,4,6),\; (4,5,7),\; (5,6,1),\;
  (6,7,2),\; (7,1,3),
\end{equation}
encoding the rule: for each triple $(a,b,c)$,
$e_a \cdot e_b = +e_c$, with cyclic permutations carrying a positive
sign and anti-cyclic permutations carrying a negative sign
\cite{Baez2002}.  Every imaginary unit squares to $-e_0$.  Among
\emph{unital} real algebras, $\OO$ is the unique normed division
algebra of dimension~8 (Hurwitz); relaxing ``division algebra'' to
the absence of zero divisors (equivalently, bijectivity of left and
right multiplication by every nonzero element) admits further
examples, including the non-unital Okubo algebra.

The \emph{norm form} is $\Nrm(x) = x \bar{x} = \sum_{i=0}^7 x_i^2$
for $x = \sum_i x_i e_i$, where $\bar{x} = x_0 e_0 - \sum_{i=1}^7
x_i e_i$ is the conjugate.  The composition property
$\Nrm(xy) = \Nrm(x)\Nrm(y)$ holds.

\subsection{The $E_8$ lattice}

In Wilson's setup~\cite{Wilson2009}, the $E_8$ lattice is
$L = D_8^+$, the set of vectors in~$\RR^8$ whose coordinates (in the
basis $\{e_0, \ldots, e_7\}$) are either all integers with even
coordinate sum, or all half-integers with odd coordinate sum
\cite{ConwaySmith2003}.  It has
240~roots (vectors of squared norm~2).  Under the standard octonion
multiplication, $L$~is closed: $L \cdot L \subseteq L$
(Lemma~\ref{lem:L-order}, certified in Table~\ref{tbl:appendix-LL}).
At this normalization $e_0 \notin L$ (the lattice is even), so $L$~is
a non-unital order; as a lattice it is a $\sqrt{2}$-scaled copy of a
maximal order of the integral octonions, the unit-scale octavians,
for which closure and maximality hold in the classical unital sense
\cite{Coxeter1946, ConwaySmith2003}.  For a modern survey of the
integral-octonion history, see~\cite{Petersson2018}.

\subsection{Wilson's Leech lattice}\label{subsec:wilson-leech}

Define $s = \tfrac{1}{2}(-e_0 + e_1 + \cdots + e_7)$ and
$\bar{s} = \tfrac{1}{2}(-e_0 - e_1 - \cdots - e_7)$, both of
norm~$\Nrm(s) = \Nrm(\bar{s}) = 2$.  Then $s \in L$ (it is a
root in the half-integer branch, with an odd coordinate sum), while $\bar{s} \notin L$
(its coordinate sum is~$-4$, which falls in the excluded even
branch of the half-integer case).  Let
\begin{equation}\label{eq:disp1}
  L\bar{s} = \{ \ell \cdot \bar{s} : \ell \in L \}, \qquad
  Ls = \{ \ell \cdot s : \ell \in L \}.
\end{equation}
Both $L\bar{s}$ and~$Ls$ are sublattices of~$L$,\footnote{$L\bar
s \subseteq L$ even though $\bar s \notin L$: each $\ell \cdot
\bar s$ with $\ell$ a root of~$L$ lies in~$L$ by direct
computation, and the inclusion extends by $\ZZ$-linearity in the
left factor. $L\bar s$ is therefore a $\ZZ$-subgroup of~$L$, and
is moreover closed under left-multiplication by~$L$ (Wilson's
condition~(2) for the untwisted product, implicit in the sublattice
argument). For~$Ls$, the subgroup
property is likewise immediate from $s \in L$.  Unlike~$L\bar s$,
the sublattice~$Ls$ is not closed under the octonion product
itself ($Ls \cdot Ls \not\subseteq Ls$).  That failure is
precisely the obstruction resolved in this paper
(Remark~\ref{rem:Ls-not-closed}).}
with
$[L : L\bar{s}] = [L : Ls] = 16$,
$L\bar{s} + Ls = L$, and $L\bar{s} \cap Ls = 2L$.

\begin{definition}[Wilson]\label{def:wilson-leech}
The \emph{Leech lattice} is
\begin{equation}\label{eq:disp2}
  \LL = \bigl\{ (x,y,z) \in L^3 :
    \text{(1) } x,y,z \in L;\;
    \text{(2) } x{+}y,\, x{+}z,\, y{+}z \in L\bar{s};\;
    \text{(3) } x{+}y{+}z \in Ls
  \bigr\}.
\end{equation}
\end{definition}

The 196{,}560 minimal vectors of~$\LL$ decompose into three families:

\smallskip
\begin{tabular}{@{}lll@{}}
\toprule
Family & Formula & Count \\
\midrule
Type~1 & $(2\lambda, 0, 0)$ and cyclic permutations & $3 \times 240 = 720$ \\
Type~2 & $(\lambda\bar{s},\; (\lambda\bar{s})j,\; 0)$ and cyclic
  permutations & $3 \times 240 \times 16 = 11{,}520$ \\
Type~3 & $((\lambda s)j,\; \lambda k,\; (\lambda j)k)$ and cyclic
  permutations & $3 \times 240 \times 16^2 = 184{,}320$ \\
\bottomrule
\end{tabular}
\smallskip

\noindent
Here $\lambda$ ranges over the 240~roots of~$L$, and
$j, k \in J = \{\pm e_t : 0 \le t \le 7\}$ (16~unit octonions).

% ===============================================================
\section{The construction}\label{sec:construction}
% ===============================================================

\begin{definition}[Transposition]\label{def:twist}
Let $\sigma = (i\; j)$ be a transposition on $\{1, \ldots, 7\}$,
extended to the linear involution $\sigma \colon \RR^8 \to \RR^8$
that fixes $e_0$ and swaps $e_i \leftrightarrow e_j$.
\end{definition}

\begin{definition}[$\sigma$-twisted octonion product]\label{def:twisted-product}
The \emph{$\sigma$-twisted octonion product}~$\cdot_\sigma$ on~$\OO$
is the $\sigma$-twisted variant of the standard octonion
product~\eqref{eq:fano}:
\begin{equation}\label{eq:twist-formula}
  x \cdot_\sigma y \;:=\; \sigma\bigl(\sigma(x) \cdot \sigma(y)\bigr).
\end{equation}
\end{definition}

This is another octonion product on $\OO$: by definition,
$\sigma \colon (\OO, \cdot) \to (\OO, \cdot_\sigma)$ is an algebra
isomorphism, $\sigma(x \cdot y) = \sigma(x) \cdot_\sigma \sigma(y)$.

\begin{definition}[Triple product with $\ZZ_3$ routing]
\label{def:triple}
Decompose $\RR^{24} = \OO_1 \oplus \OO_2 \oplus \OO_3$ into three
copies of~$\RR^8$ (blocks indexed $\alpha = 1, 2, 3$).  For
$u = (x, y, z)$ and $v = (x', y', z')$ with $x, x' \in \OO_1$, etc.,
define the \emph{triple product}
\begin{equation}\label{eq:triple}
  u \star v = (P_1,\, P_2,\, P_3),
\end{equation}
where each block $P_\alpha$ is the sum of three $\sigma$-twisted
octonion products:
\begin{align}
  P_1 &= x \cdot_\sigma x' + z \cdot_\sigma y' + y \cdot_\sigma z',
  \label{eq:P1} \\
  P_2 &= y \cdot_\sigma y' + x \cdot_\sigma z' + z \cdot_\sigma x',
  \label{eq:P2} \\
  P_3 &= z \cdot_\sigma z' + y \cdot_\sigma x' + x \cdot_\sigma y'.
  \label{eq:P3}
\end{align}
The routing rule is: a product of blocks $\alpha$ and~$\beta$ lands in
block~$\gamma$, where $\{\alpha, \beta, \gamma\} = \{1,2,3\}$ for
$\alpha \ne \beta$, and in block~$\alpha$ when $\alpha = \beta$.
All nine block-pair products use the same multiplication~$\cdot_\sigma$.
\end{definition}

% ===============================================================
\section{Proof of closure}\label{sec:proof}
% ===============================================================

We prove Theorem~\ref{thm:main} by showing that the $\sigma$-twisted
octonion product preserves each of Wilson's three conditions
(Definition~\ref{def:wilson-leech}).  Throughout, and without loss of
generality, $\sigma = (1\;2)$.\footnote{The 21 transpositions of
imaginary basis units form a single orbit under conjugation by the
order-21 group of coordinate permutations preserving the oriented Fano
triples~\eqref{eq:fano}.  Each such permutation is an octonion
automorphism that fixes $s$ and~$\bar s$ and preserves~$L$, so
conjugating by it carries the proof below for $(1\;2)$ to any other
transposition.}  From
the given choice of octonion basis and Fano triples, the
load-bearing effect of $\sigma$ is that it \emph{moves} the
sublattices $L\bar s$ and~$Ls$, while fixing~$L$ setwise
(Lemma~\ref{lem:sigma-L}) and carrying the octonion product to an
isomorphic one (Definition~\ref{def:twisted-product}).
Lemmas~\ref{lem:L-sigmaLsbar} and~\ref{lem:sigmaLs-closed}, the
closure inclusions for the moved sublattices, are where the work of
the proof lives.

\subsection{The four lemmas}

\begin{lemma}\label{lem:sigma-L}
$\sigma(L) = L$.
\end{lemma}

\begin{proof}
The lattice $L = D_8^+$ consists of all vectors in~$\RR^8$ whose
coordinates are either all integers with even coordinate sum, or all
half-integers with odd coordinate sum (Section~\ref{sec:prelim}).  The
transposition~$\sigma$ permutes two coordinates, which preserves both
the integer/half-integer parity and the coordinate sum.  Hence
$\sigma(L) \subseteq L$.  Since $\sigma$ is an involution,
$\sigma(L) = L$.
\end{proof}

For each of the next three lemmas, $L$, $\sigma(L\bar s)$, and
$\sigma(Ls)$ are full-rank $\ZZ$-lattices of rank~$8$ in~$\RR^8$.  By
$\ZZ$-bilinearity of the octonion product, an inclusion $A \cdot B
\subseteq C$ between such lattices holds if and only if it holds on a
chosen pair of $\ZZ$-bases: if $\{a_i\}_{i=1}^8$ generates $A$ and
$\{b_j\}_{j=1}^8$ generates $B$, then $A\cdot B \subseteq C$ is
equivalent to each of the 64 products $a_i\cdot b_j$ lying in
$\ZZ\langle c_1,\ldots,c_8\rangle$ for a chosen $\ZZ$-basis
$\{c_k\}_{k=1}^8$ of~$C$.  (The forward direction is the substantive
one and uses bilinearity; the reverse direction is immediate
since each $a_i\in A$ and each $b_j\in B$.)  Each lemma below is
proved by exhibiting that finite list of integer coefficients
explicitly.  Concrete
$\ZZ$-bases for $L$, $\sigma(Ls)$, and $\sigma(L\bar s)$ are fixed in
Appendix~\ref{sec:appendix-tables} (denoted there $L_k$, $M_k$, $N_k$
respectively); the 64-entry coefficient tables for the three
lemmas appear as Tables~\ref{tbl:appendix-LL},
\ref{tbl:appendix-LN}, and~\ref{tbl:appendix-MM}.

\begin{lemma}\label{lem:L-order}
$L \cdot L \subseteq L$.
\end{lemma}

\begin{proof}
The 64 products $L_i \cdot L_j$ are integer combinations of
$\{L_1,\ldots,L_8\}$; the explicit coefficients are
Table~\ref{tbl:appendix-LL}.  By $\ZZ$-bilinearity, $L \cdot L
\subseteq L$.  This is the norm-2-scale form of the classical fact
that the maximal orders of the integral octonions are closed under
multiplication \cite{Coxeter1946, ConwaySmith2003}.
\end{proof}

\begin{lemma}\label{lem:L-sigmaLsbar}
$L \cdot \sigma(L\bar{s}) \subseteq \sigma(L\bar{s})$.
\end{lemma}

\begin{proof}
The 64 products $L_i \cdot N_j$ are integer combinations of
$\{N_1,\ldots,N_8\}$; the explicit coefficients are
Table~\ref{tbl:appendix-LN}.  By $\ZZ$-bilinearity,
$L \cdot \sigma(L\bar s) \subseteq \sigma(L\bar s)$.
\end{proof}

\begin{lemma}\label{lem:sigmaLs-closed}
$\sigma(Ls) \cdot \sigma(Ls) \subseteq \sigma(Ls)$.
\end{lemma}

\begin{proof}
The 64 products $M_i \cdot M_j$ are integer combinations of
$\{M_1,\ldots,M_8\}$; the explicit coefficients are
Table~\ref{tbl:appendix-MM}.  By $\ZZ$-bilinearity,
$\sigma(Ls) \cdot \sigma(Ls) \subseteq \sigma(Ls)$.
\end{proof}

\begin{remark}\label{rem:Ls-not-closed}
The sublattice~$Ls$ is \emph{not} closed under the standard octonion
product, and $\sigma(Ls) \ne Ls$: with $\sigma = (1\;2)$, the
element $v = (-1, 0, 1, 0, 0, 1, 0, 1)$ lies in~$\sigma(Ls)
\setminus Ls$.  The twist therefore moves the block-sum sublattice
off~$Ls$ onto the closed sublattice~$\sigma(Ls)$
(Lemma~\ref{lem:sigmaLs-closed}).  Likewise $\sigma(L\bar s) \ne
L\bar s$: $w = (0, 0, -1, 0, -1, 0, -1, -1) \in \sigma(L\bar s)
\setminus L\bar s$.
\end{remark}

\subsection{Blockwise form of the triple product}

Since each $\sigma$-twisted product has the form
$a \cdot_\sigma b = \sigma(\sigma(a) \cdot \sigma(b))$ and $\sigma$
is linear, each block $P_\alpha$ is $\sigma$ applied to a sum of
standard octonion products of $\sigma$-images of the inputs.
Write $X = \sigma(x)$, $Y = \sigma(y)$, $Z = \sigma(z)$,
$X' = \sigma(x')$, $Y' = \sigma(y')$, $Z' = \sigma(z')$:
\begin{align}
  P_1 &= \sigma\bigl(X \cdot X' + Z \cdot Y' + Y \cdot Z'\bigr),
  \label{eq:P1-std} \\
  P_2 &= \sigma\bigl(Y \cdot Y' + X \cdot Z' + Z \cdot X'\bigr),
  \label{eq:P2-std} \\
  P_3 &= \sigma\bigl(Z \cdot Z' + Y \cdot X' + X \cdot Y'\bigr).
  \label{eq:P3-std}
\end{align}
Taking sums:
\begin{itemize}
\item \textbf{Block sum.}
  $P_1 + P_2 + P_3 = \sigma\bigl((X+Y+Z) \cdot (X'+Y'+Z')\bigr)$.
\item \textbf{Pairwise sums.}  For example,
  $P_1 + P_2 = \sigma\bigl(X \cdot (X'+Z') + Y \cdot (Y'+Z') +
  Z \cdot (X'+Y')\bigr)$, and similarly for the other two pairs.
\end{itemize}

\subsection{Wilson's condition~(1): each block lies in~$L$}

\begin{proposition}\label{prop:cond1}
For all $u, v \in \LL$, each block $P_\alpha$ of $u \star v$ lies
in~$L$.
\end{proposition}

\begin{proof}
Consider $P_1 = \sigma(X \cdot X' + Z \cdot Y' + Y \cdot Z')$.
Since $u \in \LL$, condition~(1) gives $x \in L$, so
$X = \sigma(x) \in L$ by Lemma~\ref{lem:sigma-L}.  Similarly
$X', Y, Y', Z, Z' \in L$.
By Lemma~\ref{lem:L-order}, each product $X \cdot X'$, $Z \cdot Y'$,
$Y \cdot Z'$ lies in~$L$.  Their sum lies in~$L$ (closed under
addition).  Applying $\sigma$ and using Lemma~\ref{lem:sigma-L},
$P_1 \in L$.  The same argument applies to
$P_2$ and~$P_3$.
\end{proof}

\subsection{Condition~(2): pairwise sums lie in~$L\bar{s}$}

\begin{proposition}\label{prop:cond2}
For all $u, v \in \LL$, the pairwise sums $P_\alpha + P_\beta$ lie
in~$L\bar{s}$.
\end{proposition}

\begin{proof}
We show $P_1 + P_2 \in L\bar{s}$; the other two cases follow by the
$\ZZ_3$ symmetry of the routing.

From \eqref{eq:P1-std}--\eqref{eq:P2-std}:
\begin{equation}\label{eq:disp3}
  P_1 + P_2 = \sigma\bigl(X \cdot (X' + Z') + Y \cdot (Y' + Z')
    + Z \cdot (X' + Y')\bigr).
\end{equation}

Since $v \in \LL$, condition~(2) gives $x' + z' \in L\bar{s}$, so
$X' + Z' = \sigma(x' + z') \in \sigma(L\bar{s})$.  Similarly
$Y' + Z' \in \sigma(L\bar{s})$ and $X' + Y' \in \sigma(L\bar{s})$.
Since $u \in \LL$, we have $X, Y, Z \in L$ (by Lemma~\ref{lem:sigma-L}).

By Lemma~\ref{lem:L-sigmaLsbar}, each product $X \cdot (X'+Z')$, $Y \cdot (Y'+Z')$,
$Z \cdot (X'+Y')$ lies in~$\sigma(L\bar{s})$.  Their sum lies
in~$\sigma(L\bar{s})$.  Applying the outer~$\sigma$:
\begin{equation}\label{eq:disp4}
  P_1 + P_2 = \sigma(\text{element of } \sigma(L\bar{s}))
  \in \sigma(\sigma(L\bar{s})) = L\bar{s},
\end{equation}
since $\sigma^2 = \mathrm{id}$.
\end{proof}

\subsection{Condition~(3): the block sum lies in~$Ls$}

This is the crux of the proof.  For the standard (untwisted) product
the block sum would require $Ls \cdot Ls \subseteq Ls$, which does
not hold (Remark~\ref{rem:Ls-not-closed}).  The twist redirects the
computation through~$\sigma(Ls)$, where closure holds.

\begin{proposition}\label{prop:cond3}
For all $u, v \in \LL$, the block sum $P_1 + P_2 + P_3$ lies in~$Ls$.
\end{proposition}

\begin{proof}
From \eqref{eq:P1-std}--\eqref{eq:P3-std}:
\begin{equation}\label{eq:disp5}
  P_1 + P_2 + P_3 = \sigma\bigl((X+Y+Z) \cdot (X'+Y'+Z')\bigr).
\end{equation}

Since $u \in \LL$, condition~(3) gives $x + y + z \in Ls$, so
\begin{equation}\label{eq:disp6}
  X + Y + Z = \sigma(x+y+z) \in \sigma(Ls).
\end{equation}
Similarly $X' + Y' + Z' \in \sigma(Ls)$.

By Lemma~\ref{lem:sigmaLs-closed}, $\sigma(Ls) \cdot \sigma(Ls) \subseteq \sigma(Ls)$, so
\begin{equation}\label{eq:disp7}
  (X+Y+Z) \cdot (X'+Y'+Z') \in \sigma(Ls).
\end{equation}
Applying the outer~$\sigma$:
\begin{equation}\label{eq:disp8}
  P_1 + P_2 + P_3 \in \sigma(\sigma(Ls)) = Ls. \qedhere
\end{equation}
\end{proof}

This completes the proof of Theorem~\ref{thm:main}.


% ===============================================================
\section{Algebraic properties}\label{sec:properties}
% ===============================================================

Throughout this section, $\sigma = (1\;2)$, as in
Section~\ref{sec:proof}.

\subsection{Identities}\label{sec:identities}

The order $(\LL, +, \star)$ has been tested for several standard
algebraic identities.  Samples are drawn uniformly (by family weight)
from $\operatorname{Min}(\LL) \times \operatorname{Min}(\LL)$ for the
binary identities, from $\operatorname{Min}(\LL)$ for the cube and
quartic identities, and from $\operatorname{Min}(\LL)^3$ for the
symmetric-composition identity.  Sample size is $N = 10^6$ per property; exact
integer arithmetic
throughout.\footnote{\scriptref{verify\_section5\_properties.py}.}  The
standard error is $\mathrm{SE} =
\sqrt{p(1-p)/N}$, where $p$ is the fraction of samples on which the
identity holds.

\begin{center}
\begin{tabular}{@{}lll@{}}
\toprule
Property & Holds? & Fraction of samples (\% $\pm$ SE) \\
\midrule
Multiplicative identity & No & n/a (structural) \\
Commutativity $u\star v = v\star u$ & No &
  $(0.018 \pm 0.0014)\%$ \\
Norm multiplicativity $\Nrm(u\star v) = \Nrm(u)\Nrm(v)$ & No &
  $(46.77 \pm 0.050)\%$ \\
Left alternativity $(x\star x)\star y = x\star(x\star y)$ & No &
  $(0.188 \pm 0.0043)\%$ \\
Right alternativity $(x\star y)\star y = x\star(y\star y)$ & No &
  $(0.188 \pm 0.0043)\%$ \\
Flexibility $(x\star y)\star x = x\star(y\star x)$ & No &
  $(8.37 \pm 0.028)\%$ \\
Cube power-associativity $x^2\star x = x\star x^2$ & No &
  $(30.70 \pm 0.046)\%$ \\
Quartic power-associativity\footnotemark{} & No &
  $(12.75 \pm 0.033)\%$ \\
Symmetric composition $\langle u\star v, w\rangle =
  \langle u, v\star w\rangle$ & No &
  $(25.37 \pm 0.044)\%$ \\
\bottomrule
\end{tabular}
\end{center}
\footnotetext{All five degree-4 parenthesizations of $x \star x \star
x \star x$ coincide.}

The failure of norm multiplicativity rules out $(\LL, +, \star)$
being a composition algebra; the failure of symmetric composition
rules out, independently, its being a symmetric composition algebra
in the sense of Okubo and para-Hurwitz
(Alberto Elduque \cite{Elduque1996, Elduque2023_IsotropicNorm};
Susumu Okubo \cite{Okubo1995_Book}).  Cube and quartic
power-associativity are reported separately; the quartic identity was
strictly stronger on every sample: in the joint exact-arithmetic
count over the same $10^6$ samples, no quartic-passing sample failed
the cube identity, while the cube-passing set was roughly
$2.4\times$ larger.

\begin{remark}\label{rem:product-norms}
Exhaustively over the minimal shell, the self-product norms
$\Nrm(u \star u)$ of the 196{,}560 minimal vectors take values in
$\{16, 32, 48, 64, 80, 96, 128\}$: neither $0$ nor $8$ occurs, so the
minimal shell has no idempotent and no square-zero
vector.\footnote{\scriptref{verify\_idempotents\_min\_shell.py}.}  On
$10^6$ sampled pairs, the pair-product norms
$\Nrm(u \star v)$ took only the values
$\{16, 32, 48, 64, 80, 96, 112, 128\}$, all multiples of~16, with
mode~64 (the norm-multiplicative value $\Nrm(u)\Nrm(v) = 8 \times
8$).\footnote{\scriptref{verify\_pair\_norm\_histogram.py}.}  Both shell questions are settled on all of~$\LL$ in
Sections~\ref{sec:idempotents} and~\ref{sec:squarezero}: the order has
no idempotent at all; it does, however, contain square-zero vectors,
the smallest at norm~12.
\end{remark}

\begin{remark}
The ambient algebra $(\RR^{24}, +, \star)$ has no multiplicative
identity: any putative left identity~$(a,b,c)$ would have to act on
an arbitrary~$(x',y',z')$ by the block matrix
$\bigl(\begin{smallmatrix} a & c & b \\ c & b & a \\ b & a & c
\end{smallmatrix}\bigr)$ (entries acting by left-multiplication
in~$\OO$), and no triple $(a,b,c)$ makes this the identity on all
three blocks simultaneously.  In particular, $(e_0, e_0, e_0)$ is
not an identity: it maps $(x', y', z')$ to $(x'+y'+z',\, x'+y'+z',\,
x'+y'+z')$.  Consequently $(\LL, +, \star)$ is a non-unital order.
\end{remark}

The order does not satisfy any of the classical algebraic identities
(associativity, alternativity, flexibility, power-associativity,
commutativity).  The defining property is \emph{closure}: the Leech
lattice is stable under the product.

\subsection{Idempotents}\label{sec:idempotents}

The ambient algebra $(\RR^{24}, +, \star)$ has exactly eight
idempotents, a symbolic classification carried out and verified in
exact arithmetic.\footnote{\scriptref{verify\_idempotent\_classification.py}.}
Writing $\varepsilon_1 = (e_0, 0, 0)$, $\varepsilon_2 = (0, e_0, 0)$,
$\varepsilon_3 = (0, 0, e_0)$ and $\omega = \varepsilon_1 +
\varepsilon_2 + \varepsilon_3$, they are
\begin{equation}\label{eq:idempotents}
  0, \qquad \varepsilon_1, \ \varepsilon_2, \ \varepsilon_3, \qquad
  \tfrac13 \omega, \qquad
  \varepsilon_i - \tfrac13 \omega \ \ (i = 1, 2, 3).
\end{equation}
Note that every block of every one of them is a multiple of~$e_0$.

Each nonzero idempotent spans a rational ray that meets~$\LL$, and the
least positive multiple that does
is\footnote{\scriptref{verify\_idempotent\_lattice\_rescaling.py}.}
\begin{equation}\label{eq:idempotent-multiples}
  4\varepsilon_i \ (\Nrm = 16), \qquad
  4\omega \ (\Nrm = 48), \qquad
  6\bigl(\varepsilon_i - \tfrac13\omega\bigr) \ (\Nrm = 24).
\end{equation}

\subsection{The automorphism group}\label{sec:autgroup}

$\Aut(\LL, +, \star)$ is finite, of order~$36$, with structure
$C_6 \times S_3$.  Every automorphism is a blockwise octonion
automorphism followed by a permutation of the three blocks; in the
explicit demonstration below, blockwise means with respect to the
twisted algebra $(\OO, \cdot_\sigma)$ of
Definition~\ref{def:twisted-product}.  The~$S_3$
is the block permutation of Definition~\ref{def:triple}, and the~$C_6$
is generated by the signed permutation of imaginary units
$e_1 \mapsto -e_5$, $e_2 \mapsto -e_7$, $e_3 \mapsto e_2$,
$e_4 \mapsto -e_4$, $e_5 \mapsto e_6$, $e_6 \mapsto e_1$,
$e_7 \mapsto -e_3$, with $e_0$ fixed.

The containment in the Conway group can already be settled at the level
of the ambient algebra, before the lattice is imposed.  As a real
algebra, $\RR^{24} = \OO^3$
splits as $D \oplus T$, where $D = \{(a,a,a) : a \in \OO\}$ (of
dimension~8) and $T = \{(p,q,r) \in \OO^3 : p+q+r = 0\}$ (of
dimension~16) are two-sided ideals with $D \star T = T \star D = 0$;
every
$\star$-automorphism preserves them and acts on each by an octonion
automorphism, composed with a permutation of the blocks, so that
$\Aut(\RR^{24}, \star) = G_2 \times G_2 \times S_3$.  An automorphism of
the \emph{real} octonions preserves the positive-definite octonion norm,
and $D \perp T$; hence every $\star$-automorphism of $\RR^{24}$ is
orthogonal, and
\begin{equation}\label{eq:disp9}
  \Aut(\LL, +, \star) \;\subsetneq\; \Co_0 .
\end{equation}
The inclusion is strict because $-I_{24} \in \Co_0$ does not
preserve~$\star$.  The derivation, the complete enumerations, and an
independent verification are recorded separately in a companion
note.\footnote{\cite{repo}~%
\href{https://github.com/koeplinger/leech\_alg/blob/main/paper/automorphism\_group\_2026-07-12.pdf}%
{\texttt{paper/automorphism\_group\_2026-07-12.pdf}};
\scriptref{verify\_aut\_lambda\_star.py}
\gapref{aut\_lambda\_star.g}.}

\subsection{Square-zero elements}\label{sec:squarezero}

In the ambient algebra, $u = (x, y, z)$ satisfies $u \star u = 0$ if
and only if $x$, $y$, $z$ are purely imaginary, of equal norm, and sum
to zero, again a symbolic classification.  Three vectors of equal length summing to zero lie at $120^\circ$
to one another; with their negatives they are the six vertices of a
regular hexagon, and $x$, $y$, $z$ are three alternating vertices, a
\emph{hexagonal triple} in $\Im(\OO) = \RR^7$.  The solution set is
a cone of dimension~12; it lies in~$T$ and meets~$D$ only
in~$0$.\footnote{\scriptref{verify\_square\_zero\_classification.py}.}

The lattice does contain such elements.  A vector $(x, y, z) \in \LL$
squares to zero exactly when $x$, $y$, $z$ lie in $L\bar s \cap
\Im(\OO)$, have equal norms and sum to zero.  The 126 vectors of norm~4
in $L\bar s \cap \Im(\OO)$ form an $E_7$ root system.  Each of its 126
roots has exactly 32 roots at $120^\circ$; the third block is then
forced, so the ordered hexagonal triples number $126 \times 32 = 4{,}032$,
one square-zero vector of~$\LL$ each, all of norm~12.  Every square-zero
norm of~$\LL$ lies in $12\ZZ$ ($\Nrm(u) = 3\Nrm(x)$ with $\Nrm(x)$ even,
and every norm of~$\LL$ lies in $4\ZZ$), and the $4{,}032$ root triples
are exactly the norm-12 stratum; higher strata are realized, for
example, the blockwise sum of two blockwise orthogonal root triples is
square-zero of norm~24.\footnote{\scriptref{verify\_square\_zero\_norm24\_witness.py}.}

\subsection{Span of the image}\label{sec:span}

Let $S := \ZZ\text{-span}\{u \star v : u, v \in \LL\}$.  By bilinearity, $S$ is
generated by the $24 \times 24 = 576$ basis-pair products of a
$\ZZ$-basis of~$\LL$, so its index in~$\LL$ is exactly
computable:\footnote{\scriptref{verify\_product\_span\_index.py}
\scriptref{verify\_product\_span\_structure.py}.}
\begin{equation}\label{eq:span-index}
  [\LL : S] \;=\; 65{,}536 \;=\; 2^{16}.
\end{equation}
The index counts the cosets of~$S$ in~$\LL$: the lattice~$\LL$
decomposes into $[\LL : S]$ disjoint translates of~$S$, and every
product, together with every integer combination of products, lies
in a single one of them.  A span-surjective
product would have index~1.  Here, $2^{16} - 1$ of every $2^{16}$
cosets are unreachable.  Moreover $2\LL \subseteq S$, so $S$ is the
full preimage of an 8-dimensional $\FF_2$-subspace of $\LL/2\LL$,
and $\LL/S \cong (\ZZ/2)^{16}$.  The defect is entirely 2-primary,
at the same prime that governs the mod-2 structure of
Section~\ref{app:mod2quotient}.  Identifying the subspace $S/2\LL$
inside $\LL/2\LL$ structurally is open.

\subsection{Closure of permutation cycles}\label{sec:census}

The transposition~$\sigma$ is a special case of the permutations of the
seven imaginary octonion basis elements, described by the symmetric
group~$S_7$.  For a general such permutation~$\pi$, the product
$x \cdot_\pi y := \pi(\pi(x) \cdot \pi(y))$ may or may not close
on~$\LL$.  Testing closure exactly on the 576 basis-pair products of a
$\ZZ$-basis of~$\LL$, we find for every cycle type of~$S_7$:\footnote{\scriptref{verify\_all\_permutations\_exact.py}.}
\begin{center}
\begin{tabular}{@{}lrrr@{}}
\toprule
Cycle type of $\pi$ & Permutations & Close & Rate \\
\midrule
$2$-cycle     & $21$  & $21$  & $100\%$ \\
$3$-cycle     & $70$  & $35$  & $50\%$ \\
$(2,2)$       & $105$ & $42$  & $40\%$ \\
$4$-cycle     & $210$ & $0$   & $0\%$ \\
$(3,2)$       & $420$ & $126$ & $30\%$ \\
$5$-cycle     & $504$ & $252$ & $50\%$ \\
$(2,2,2)$     & $105$ & $21$  & $20\%$ \\
$(3,3)$       & $280$ & $112$ & $40\%$ \\
$(4,2)$       & $630$ & $294$ & $46.7\%$ \\
$6$-cycle     & $840$ & $210$ & $25\%$ \\
$(3,2,2)$     & $210$ & $105$ & $50\%$ \\
$(4,3)$       & $420$ & $168$ & $40\%$ \\
$(5,2)$       & $504$ & $42$  & $8.3\%$ \\
$7$-cycle     & $720$ & $336$ & $46.7\%$ \\
\midrule
Total         & $5{,}039$ & $1{,}764$ & $35.0\%$ \\
\bottomrule
\end{tabular}
\end{center}
The identity (the untwisted product) does not close.  The $4$-cycles
are the only type failing without exception; the $(5,2)$ type has the
lowest nonzero rate.  The $3$-cycles, $5$-cycles and $(3,2,2)$ types
close at exactly half.  For the $3$-cycles the closing half admits a
description: for each three-element subset of imaginary indices, one
of the two cyclic orientations closes.  No such description is known
for the other cycle types, and none is known for the $0\%$ of the
$4$-cycles.

% ===============================================================
\section{Related work}\label{sec:related}
% ===============================================================

\textbf{Wilson (2009)} \cite{Wilson2009} characterizes~$\LL$ as
$\{(x,y,z) \in L^3 : \text{conditions 1--3}\}$, gives an explicit
listing of the 196{,}560 minimal vectors in the form
$3 \times 240 \times (1 + 16 + 16 \times 16)$, and proves that~$\LL$
is even and self-dual.  The sublattice framework is
the setting used in the present paper.

\textbf{Dixon (2010)} \cite{Dixon2010} constructs~$\LL$ using the
octonion XY-product\footnote{The construction has an earlier precursor
in Dixon~\cite{Dixon1995}, where a closely related octonionic
description of~$\LL$ is proposed.  A 2005 preprint of
Dixon's~\cite{Dixon2005} carried a formula whose third
family of minimal vectors required a small correction, which Wilson
\cite{Wilson2009} supplied in a bracketed note while independently
deriving the sublattice characterization.  The situation is similar
to the Kirmse--Bruck--Coxeter relationship recorded in
Appendix~\ref{app:history}: an early, structurally correct proposal
needing correction.}
\begin{equation}\label{eq:disp10}
  A \circ_{XY} B \;=\; (A X)(Y^{\dagger} B),
\end{equation}
where $X, Y \in \OO$ are chosen so that $Y X^{\dagger}$ is a unit
octonion (the identity of the new multiplication).  When $Y^{\dagger} = X^{-1}$
the XY-product specializes to the X-product $A \circ_X B = (AX)(X^{-1}
B)$; Dixon uses this X-product with $X = \ell_0 = \frac{1}{2}(1 + e_1
+ \cdots + e_7)$ to obtain~$E_8$, and general XY-products for
higher-dimensional laminated lattices including~$\LL$.  We have
independently reproduced this construction in the supplemental
code~\cite{repo}.

\textbf{Baez--Egan (2014)} \cite{Baez2014} study closure
properties of integer sublattices of the exceptional Jordan algebra
$\mathfrak{h}_3(\OO)$.  In parts~9--10 of the series, Baez considers
the 27-dimensional lattice $\ZZ^3 \oplus L_L \subset \mathfrak{h}_3(\OO)$,
with integers on the diagonal and an off-diagonal Leech sublattice
$L_L \cong \LL$ embedded as an octonionic triple.  This lattice is
\emph{not} closed under the Jordan product $x \circ y = xy + yx$,
but \emph{is} closed under the doubled product $2(xy + yx)$.
Egan further enumerated at least $17{,}280$ distinct off-diagonal
Leech embeddings $L_L \subset \OO^3$ for which the associated
27-dimensional lattice is closed under this doubled product
(see part~10 of~\cite{Baez2014} for Egan's computational counts).
Baez left open the question whether any such lattice is closed under
the undoubled Jordan product, and conjectures that the answer is
negative.  We
have independently reproduced the full Baez--Egan verification chain:
the $E_8$ simple-roots data, the lemma $L_1 \cap L_2 = 2L_0$, the
Leech-lattice theorem, the doubled-Jordan-product closure, the
derivation of the lower bound $244{,}035{,}421$ on the number of
$E_8^3$-orbit Leech embeddings,\footnote{An orbit-stabilizer bound, $\lceil 3!\,\lvert W(E_8)\rvert^3 / \lvert\Co_0\rvert\rceil$; \scriptref{verify\_egan\_lower\_bound.py}.} and Egan's count of at least
$17{,}280$ Jordan rings.  This last factors as $17{,}280 = 270
\times 64$: the $270$ is the number of sublattices $L_1 \subset L$
isometric to $\sqrt{2}\,L$ (part~11 of \cite{Baez2014}), equivalently
the number of $4$-dimensional maximal totally isotropic subspaces of
$L/2L \cong \FF_2^8$ under the plus-type quadratic form induced by the
$E_8$ Gram form; that mod-2 orthogonal-space picture is Tim
Silverman's (part~12 of \cite{Baez2014}).  The $64$ is
the number of $L_2$ in this collection that meet a given $L_1$ only
in $2L$, equivalently the number of $4$-dimensional complementary
isotropic subspaces.  All Baez--Egan claims hold under our reading.
Details and code are in~\cite{repo}.

\textbf{Comparison.}  The product~$\star$ is bilinear on~$\RR^{24}$
itself.  It does not require an enclosing 27-dimensional
Jordan-algebraic ambient within which the Leech sits as a subspace
(cf.~\cite{Baez2014}), and the cross-block $S_3$ extending the
$\ZZ_3$ routing is supplied by the routing of
Definition~\ref{def:triple} rather than by the octonion product
itself: Dixon's XY-product
\cite{Dixon1995, Dixon2010} is a useful device for describing each
$E_8$ block, but it is not in itself $\ZZ_3$-symmetric across the
three factors.  Whether the Leech embedding $\sigma(\LL) \subset
\OO^3$ obtained by applying $\sigma$ blockwise overlaps with one of
the embeddings catalogued by Egan, or sits outside that family, is
a question this paper does not address.

\textbf{Anatomy of the Baez--Egan closure on $\LL$.}  The
Baez--Egan closure is on the ambient 27-dimensional lattice
$\widetilde{\LL} = \ZZ^3 \oplus \LL \subset \mathfrak{h}_3(\OO)$,
not on $\LL$ alone.  For off-diagonal Hermitian matrices $X, Y \in
\mathfrak{h}_3(\OO)$ representing $u = (x,y,z)$ and $v = (x',y',z')$
in $\LL \subset \OO^3$, direct computation gives
\begin{equation}\label{eq:disp11}
  2(XY + YX) \;=\; D \;+\; R,
\end{equation}
where $R$ is an off-diagonal residue lying in $\LL$ and $D \in
\ZZ^3$ is a diagonal piece whose entries are four times the
block-pair Euclidean inner products of $u$ and $v$:
\begin{equation}\label{eq:disp12}
  D \;=\; 4\bigl(\langle x, x'\rangle + \langle y, y'\rangle,\;
                  \langle x, x'\rangle + \langle z, z'\rangle,\;
                  \langle y, y'\rangle + \langle z, z'\rangle\bigr),
\end{equation}
so $\operatorname{tr} D = 8 \langle u, v\rangle$ recovers the
$\RR^{24}$ inner product of $u$ and $v$.  Projecting onto the
off-diagonal subspace defines a $\ZZ$-bilinear closed product
$\varphi(u, v) := R$ on $\LL$: an order in the sense of
Theorem~\ref{thm:main}, but obtained by projecting from a larger
ambient with the inner-product shadow $D$ discarded.  For minimal
Leech vectors $u \in \operatorname{Min}(\LL)$, squaring gives
$\operatorname{tr} D = 8\Nrm(u) = 64$, so the projection is
non-trivial on every minimal-vector input.  The induced product
$\varphi$ is moreover commutative (the Jordan product is symmetric
in its arguments), whereas $\star$ is not
(Section~\ref{sec:properties}); the two orders are therefore
distinct bilinear products on $\RR^{24}$ that both close on $\LL$,
reaching that closure through different ambient algebras.  The
distinction is quantitative at the span level as well.  Computed by
the same Hermite-normal-form method as for~$\star$
(Section~\ref{sec:span}),\footnote{\scriptref{verify\_phi\_span\_index.py}.}
the $\ZZ$-span of $\varphi(\LL, \LL)$ has index $2^{25}$ in~$\LL$.
Unlike for~$\star$, that span does not contain~$2\LL$: its index
exceeds $[\LL : 2\LL] = 2^{24}$.

\textbf{Earlier precursors.}  Further back, several constructions
foreshadow the octonion-triple and Jordan-algebra frameworks above.
John H.~Conway and Neil J.~A.~Sloane~\cite{ConwaySloane1982} give 23
constructions of~$\LL$, one for each Niemeier lattice, including a
Turyn-style construction from three copies of~$E_8$.  James Lepowsky
and Arne Meurman~\cite{LepowskyMeurman1982} describe~$\LL$ as a
sublattice of~$E_8^3$ containing~$2 E_8^3$ and use this to generate
the Conway group~$\Co_0$ via permutations of the three $E_8$ factors
together with one extra automorphism.  Noam D.~Elkies and
Benedict H.~Gross~\cite{ElkiesGross1996}
obtain~$\LL$ from an integer form of the exceptional Jordan algebra
$\mathfrak{h}_3(\OO)$ as the orthogonal complement of a distinguished
identity element, and show that the automorphism group of this
integer form is generated by octonion reflections [\emph{octave
reflections}, in their terminology].  Each of these
constructions is additive or geometric: none defines or tests a
bilinear product on~$\LL$.

\textbf{Leonard E.~Dickson (1923), Johannes Kirmse (1924), Kurt Mahler
(1942), Richard H.~Bruck and H.\,S.\,M.~Coxeter (1946).}  For a full
historical account, see Appendix~\ref{app:history}.

\textbf{Daniele Corradetti (2026).}  A thematically adjacent recent
construction is Corradetti's order on~$E_8$ via the Okubo algebra
\cite{Corradetti2026}: a different non-classical bilinear structure
on~$E_8$, obtained by a $2$-adic conductor / saturation of an
Okubo-algebra overlattice.  The two operations involved differ from
the present mod-2 quotient (Appendix~\ref{sec:appendix-tables},
Section~\ref{app:mod2quotient}): Corradetti rescales an order and
$2$-adically saturates a sublattice, while the present paper passes
to a quotient $L \to L/2L$.  Both place $2$-primary structure at
the center of an octonion-derived order on~$E_8$.

% ===============================================================
\section{Conclusion}\label{sec:conclusion}
% ===============================================================

The Leech lattice admits an order under a $\ZZ_3$-symmetric
triple-octonion product on~$\RR^{24}$, realized in Wilson's
sublattice representation of~$\LL$.  The symbolic proof
(Section~\ref{sec:proof}) reduces closure to four lemmas, of which
the crux is Lemma~\ref{lem:sigmaLs-closed}: the sublattice
$\sigma(Ls)$, unlike~$Ls$ itself, is closed under the standard
octonion product.

Several questions remain immediately open:
\begin{itemize}
\item A structural reason for Lemma~\ref{lem:sigmaLs-closed}.  The
  proof in Appendix~\ref{sec:appendix-tables} exhibits the 64
  integer coefficients certifying $\sigma(Ls)\cdot\sigma(Ls) \subseteq
  \sigma(Ls)$, but does not explain \emph{why} $\sigma(Ls)$ is closed
  under the standard octonion product while $Ls$ is not.  The mod-2
  reduction in Section~\ref{app:mod2quotient} identifies what the
  appendix tables certify after reduction modulo~$2L$: $W :=
  \sigma(Ls)/2L$ is a subalgebra of $L/2L$, the reduction of the
  octonion order~$L$ modulo~$2$.  One candidate shape for an
  integer-level reason is ruled out: $\sigma(Ls)$ is not an ideal
  of~$L$ on either side.  Of the 64 basis-pair products, 32 of
  $L \cdot \sigma(Ls)$ and 43 of $\sigma(Ls) \cdot L$ lie outside
  $\sigma(Ls)$.\footnote{\scriptref{verify\_sigma\_Ls\_ideal\_exclusion.py}.}
  The companion sublattice $\sigma(L\bar s)$ is by contrast a
  \emph{two-sided} ideal of~$L$, of which
  Lemma~\ref{lem:L-sigmaLsbar} is the left-handed half, and the same
  asymmetry survives reduction modulo~$2L$: $V$ is an ideal
  of~$L/2L$, $W$ only a subalgebra.  An integer-level structural
  reason would replace the appendix tables by a structural argument
  resting on something other than an ideal property: an
  identification of $\sigma(Ls)$ with the image under a
  norm-preserving octonion automorphism, or with some
  module-theoretic structure over~$L$ that lifts the mod-2 picture.  (See the
  closing paragraphs of Section~\ref{app:mod2quotient} for a
  Lagrangian-polarization reframing of this open question as a
  research program.)
\item Where $\Aut(\LL, +, \star) \cong C_6 \times S_3$
  (Section~\ref{sec:autgroup}) sits inside~$\Co_0$ up to conjugacy, and
  whether the order-6 octonion automorphism generating the~$C_6$
  admits an intrinsic description.
\item Whether the order is maximal: is there a larger lattice
  $\LL' \supsetneq \LL$, containing $\LL$ with finite index, for which
  $\LL' \star \LL' \subseteq \LL'$?  Any such $\LL'$ lies inside
  $\tfrac{1}{n}\LL$ for some integer~$n$, so for each~$n$ the question
  is a finite check of the same kind used here.
\item Classification of the algebra $(\RR^{24}, +, \star)$: it is
  non-associative, non-alternative, non-flexible, and non-unital
  (as an order), placing it outside the classical families.
\item A structural characterization of which permutation cycles close.
  Every cycle type of~$S_7$ has been enumerated exactly
  (Section~\ref{sec:census}), but the pattern in the table is
  explained only for the $3$-cycles.  Open: a reason for the total
  failure of the $4$-cycles and for the exact one-half rates; and
  linear maps of~$\RR^8$ that are not coordinate permutations.
\end{itemize}

% ===============================================================
\section{Outlook}\label{sec:outlook}
% ===============================================================

A more natural classification of $(\RR^{24}, +, \star)$ may emerge
from a \emph{ternary} rather than binary viewpoint.  The Leech lattice
carries a prominent $\ZZ_3$-symmetry (the cyclic permutation of the
three octonion blocks), and~$\star$ is assembled from nine bilinear
building blocks routed by that $\ZZ_3$-action.  Ternary composition
algebras, as studied by
Elduque~\cite{Elduque1996,Elduque2000_Triality,Elduque2023_IsotropicNorm}
and others~\cite{SmithVojtechovsky2022,KamiyaOkubo2015,Okubo1995_Book},
are typically non-associative, non-commutative, non-alternative, and
non-power-associative, the same negative classification $(\RR^{24},
+, \star)$ satisfies.  Yet they carry high internal symmetry, with
order-3 automorphisms playing the central role: the Okubo algebra has
automorphism group $\mathrm{SU}(3)$, and its product is built from an
order-3 automorphism of the octonions.  Jonathan I.~Hall's
monograph~\cite{Hall2019_Moufang} establishes that Moufang loops and
groups with triality are essentially the same object.  Recasting
$\star$ as a trace-like reduction of a ternary operation compatible
with the $\ZZ_3$-block symmetry is a natural direction for further
work.

Programs that model fundamental physics on normed-division-algebra,
triality, $S_3$ symmetry, or symmetric-composition-algebra structures
\cite{GresnigtGourlayVarma2023, Koeplinger2023,
FureyHughes2025_TrioOfTrialities, MarraniCorradettiZucconi2025}
are one setting in which a non-classical bilinear product on~$\RR^{24}$
may be of interest.

% ===============================================================
\section*{Acknowledgments}
% ===============================================================

The author thanks Matthew Barley, Geoffrey M. Dixon, Holger P.
Petersson, Petr Vojt\v{e}chovsk\'{y}, and Robert A. Wilson for
valuable discussion, advice, feedback, and corrections.  The Lagrangian-decomposition
framing recorded at the end of
Section~\ref{app:mod2quotient} was prompted by feedback from Hermes
(using OpenCode with DeepSeek).

% ===============================================================
\section*{Declarations}
% ===============================================================

\textbf{Funding.}  The author received no funding for this work.

\textbf{Competing interests.}  The author has no competing interests
to declare.

\textbf{Data and code availability.}  All data and code are public:
the complete research record, including every verification script
cited in this paper, is available at the repositories of~\cite{repo}.

\textbf{Use of AI.}  AI assistance is a stated, material part of this
work; the methodology, including the role of AI and the human-run
verification protocol, is documented in
Appendix~\ref{sec:methodology}.

% ===============================================================
\appendix
\section{Explicit basis tables for Lemmas
\ref{lem:L-order}--\ref{lem:sigmaLs-closed}}\label{sec:appendix-tables}
% ===============================================================

We fix concrete $\ZZ$-bases for the lattices $L$, $\sigma(Ls)$, and
$\sigma(L\bar s)$, and tabulate the integer coefficients of the
basis-by-basis products called for in the proofs of
Lemmas~\ref{lem:L-order}, \ref{lem:L-sigmaLsbar},
and~\ref{lem:sigmaLs-closed}.  Each table has 64 entries and is
arranged so that the row labeled $(i,j)$ gives the unique 8-tuple
$(c_1,\ldots,c_8) \in \ZZ^8$ representing the corresponding product
in the basis indicated.  All entries are integer; closure under the
standard octonion product follows by $\ZZ$-bilinearity.  The same
identities have been recomputed in two independent computer-algebra
implementations as a cross-check
(Appendix~\ref{sec:methodology}).

Coordinates throughout are with respect to the basis
$\{e_0,\ldots,e_7\}$ and the standard Fano triples~\eqref{eq:fano}.

\paragraph{The basis $\{L_k\}$.} Concretely: $L_1 = s = \tfrac12(-e_0+e_1+\cdots+e_7)$, $L_2 = e_0+e_1$, $L_3 = e_0-e_1$, and $L_{k+2} = e_0+e_k$ for $k=2,\ldots,6$. The eight vectors are linearly independent over $\QQ$ and span $L$ over $\ZZ$.
\begin{align}
  L_{1} &= (\tfrac{-1}{2},\ \tfrac{1}{2},\ \tfrac{1}{2},\ \tfrac{1}{2},\ \tfrac{1}{2},\ \tfrac{1}{2},\ \tfrac{1}{2},\ \tfrac{1}{2}) \\
  L_{2} &= (1,\ 1,\ 0,\ 0,\ 0,\ 0,\ 0,\ 0) \\
  L_{3} &= (1,\ -1,\ 0,\ 0,\ 0,\ 0,\ 0,\ 0) \\
  L_{4} &= (1,\ 0,\ 1,\ 0,\ 0,\ 0,\ 0,\ 0) \\
  L_{5} &= (1,\ 0,\ 0,\ 1,\ 0,\ 0,\ 0,\ 0) \\
  L_{6} &= (1,\ 0,\ 0,\ 0,\ 1,\ 0,\ 0,\ 0) \\
  L_{7} &= (1,\ 0,\ 0,\ 0,\ 0,\ 1,\ 0,\ 0) \\
  L_{8} &= (1,\ 0,\ 0,\ 0,\ 0,\ 0,\ 1,\ 0)
\end{align}

\paragraph{The basis $\{M_k\}$.} Defined by $M_k := \sigma(L_k \cdot s)$ with $\sigma = (1\;2)$ acting on coordinates by swapping $e_1$ and $e_2$. Spans $\sigma(Ls)$ over $\ZZ$.
\begin{align}
  M_{1} &= (\tfrac{-3}{2},\ \tfrac{-1}{2},\ \tfrac{-1}{2},\ \tfrac{-1}{2},\ \tfrac{-1}{2},\ \tfrac{-1}{2},\ \tfrac{-1}{2},\ \tfrac{-1}{2}) \\
  M_{2} &= (-1,\ 0,\ 0,\ 0,\ 1,\ 0,\ 1,\ 1) \\
  M_{3} &= (0,\ 1,\ 1,\ 1,\ 0,\ 1,\ 0,\ 0) \\
  M_{4} &= (-1,\ 0,\ 1,\ 0,\ 0,\ 1,\ 0,\ 1) \\
  M_{5} &= (-1,\ 1,\ 1,\ 0,\ 0,\ 0,\ 1,\ 0) \\
  M_{6} &= (-1,\ 1,\ 0,\ 1,\ 0,\ 0,\ 0,\ 1) \\
  M_{7} &= (-1,\ 0,\ 1,\ 1,\ 1,\ 0,\ 0,\ 0) \\
  M_{8} &= (-1,\ 1,\ 0,\ 0,\ 1,\ 1,\ 0,\ 0)
\end{align}

\paragraph{The basis $\{N_k\}$.} Defined by $N_k := \sigma(L_k \cdot \bar s)$. Spans $\sigma(L\bar s)$ over $\ZZ$.
\begin{align}
  N_{1} &= (2,\ 0,\ 0,\ 0,\ 0,\ 0,\ 0,\ 0) \\
  N_{2} &= (0,\ 0,\ -1,\ 0,\ -1,\ 0,\ -1,\ -1) \\
  N_{3} &= (-1,\ -1,\ 0,\ -1,\ 0,\ -1,\ 0,\ 0) \\
  N_{4} &= (0,\ -1,\ -1,\ 0,\ 0,\ -1,\ 0,\ -1) \\
  N_{5} &= (0,\ -1,\ -1,\ -1,\ 0,\ 0,\ -1,\ 0) \\
  N_{6} &= (0,\ -1,\ 0,\ -1,\ -1,\ 0,\ 0,\ -1) \\
  N_{7} &= (0,\ 0,\ -1,\ -1,\ -1,\ -1,\ 0,\ 0) \\
  N_{8} &= (0,\ -1,\ 0,\ 0,\ -1,\ -1,\ -1,\ 0)
\end{align}

\begin{table}[H]
\centering
{\bfseries Lemma~\ref{lem:L-order}: $L \cdot L \subseteq L$.}\par\medskip
\footnotesize
\setlength{\tabcolsep}{4pt}
\begin{minipage}[t]{0.48\textwidth}
\centering
\begin{tabular}{cc|rrrrrrrr}
\toprule
$i$ & $j$ & $c_1$ & $c_2$ & $c_3$ & $c_4$ & $c_5$ & $c_6$ & $c_7$ & $c_8$ \\
\midrule
1 & 1 & -1 & -1 & -1 & 0 & 0 & 0 & 0 & 0 \\
1 & 2 & 0 & -2 & -2 & 1 & 1 & 0 & 1 & 0 \\
1 & 3 & 2 & 2 & 2 & -1 & -1 & 0 & -1 & 0 \\
1 & 4 & 0 & -2 & -2 & 0 & 1 & 1 & 0 & 1 \\
1 & 5 & 2 & 1 & 2 & -1 & -1 & 0 & 0 & -1 \\
1 & 6 & 0 & -1 & -2 & 0 & 0 & 0 & 1 & 1 \\
1 & 7 & 2 & 1 & 2 & 0 & -1 & -1 & -1 & 0 \\
1 & 8 & 2 & 2 & 2 & -1 & 0 & -1 & -1 & -1 \\
\midrule
2 & 1 & 2 & 1 & 2 & -1 & -1 & 0 & -1 & 0 \\
2 & 2 & 0 & 1 & -1 & 0 & 0 & 0 & 0 & 0 \\
2 & 3 & 0 & 1 & 1 & 0 & 0 & 0 & 0 & 0 \\
2 & 4 & 0 & 0 & -1 & 1 & 0 & 1 & 0 & 0 \\
2 & 5 & 2 & 3 & 3 & -1 & 0 & -1 & -1 & -1 \\
2 & 6 & 0 & 1 & 0 & -1 & 0 & 1 & 0 & 0 \\
2 & 7 & 0 & 0 & -1 & 0 & 0 & 0 & 1 & 1 \\
2 & 8 & 0 & 1 & 0 & 0 & 0 & 0 & -1 & 1 \\
\midrule
3 & 1 & 0 & -1 & -2 & 1 & 1 & 0 & 1 & 0 \\
3 & 2 & 0 & 1 & 1 & 0 & 0 & 0 & 0 & 0 \\
3 & 3 & 0 & -1 & 1 & 0 & 0 & 0 & 0 & 0 \\
3 & 4 & 0 & 0 & 1 & 1 & 0 & -1 & 0 & 0 \\
3 & 5 & -2 & -3 & -3 & 1 & 2 & 1 & 1 & 1 \\
3 & 6 & 0 & -1 & 0 & 1 & 0 & 1 & 0 & 0 \\
3 & 7 & 0 & 0 & 1 & 0 & 0 & 0 & 1 & -1 \\
3 & 8 & 0 & -1 & 0 & 0 & 0 & 0 & 1 & 1 \\
\midrule
4 & 1 & 2 & 2 & 2 & -1 & -1 & -1 & 0 & -1 \\
4 & 2 & 0 & 1 & 0 & 1 & 0 & -1 & 0 & 0 \\
4 & 3 & 0 & -1 & 0 & 1 & 0 & 1 & 0 & 0 \\
4 & 4 & 0 & -1 & -1 & 2 & 0 & 0 & 0 & 0 \\
4 & 5 & 0 & -1 & -1 & 1 & 1 & 0 & 1 & 0 \\
4 & 6 & 0 & 0 & -1 & 1 & 0 & 1 & 0 & 0 \\
4 & 7 & 0 & 0 & 0 & 1 & -1 & 0 & 1 & 0 \\
4 & 8 & 2 & 2 & 3 & 0 & -1 & -1 & -1 & 0 \\
\bottomrule
\end{tabular}
\end{minipage}\hfill
\begin{minipage}[t]{0.48\textwidth}
\centering
\begin{tabular}{cc|rrrrrrrr}
\toprule
$i$ & $j$ & $c_1$ & $c_2$ & $c_3$ & $c_4$ & $c_5$ & $c_6$ & $c_7$ & $c_8$ \\
\midrule
5 & 1 & 0 & -1 & -2 & 1 & 0 & 0 & 0 & 1 \\
5 & 2 & -2 & -2 & -4 & 1 & 2 & 1 & 1 & 1 \\
5 & 3 & 2 & 2 & 4 & -1 & 0 & -1 & -1 & -1 \\
5 & 4 & 0 & 0 & 0 & 1 & 1 & 0 & -1 & 0 \\
5 & 5 & 0 & -1 & -1 & 0 & 2 & 0 & 0 & 0 \\
5 & 6 & 0 & -1 & -1 & 0 & 1 & 1 & 0 & 1 \\
5 & 7 & 0 & -1 & -1 & 1 & 1 & 0 & 1 & 0 \\
5 & 8 & 0 & 0 & 0 & 0 & 1 & -1 & 0 & 1 \\
\midrule
6 & 1 & 2 & 1 & 2 & 0 & 0 & -1 & -1 & -1 \\
6 & 2 & 0 & 0 & -1 & 1 & 0 & 1 & 0 & 0 \\
6 & 3 & 0 & 0 & 1 & -1 & 0 & 1 & 0 & 0 \\
6 & 4 & 0 & -1 & 0 & 1 & 0 & 1 & 0 & 0 \\
6 & 5 & 0 & 0 & 0 & 0 & 1 & 1 & 0 & -1 \\
6 & 6 & 0 & -1 & -1 & 0 & 0 & 2 & 0 & 0 \\
6 & 7 & 2 & 2 & 3 & -1 & -1 & 0 & 0 & -1 \\
6 & 8 & 0 & -1 & -1 & 0 & 1 & 1 & 0 & 1 \\
\midrule
7 & 1 & 0 & -1 & -2 & 0 & 1 & 1 & 0 & 0 \\
7 & 2 & 0 & 1 & 0 & 0 & 0 & 0 & 1 & -1 \\
7 & 3 & 0 & -1 & 0 & 0 & 0 & 0 & 1 & 1 \\
7 & 4 & 0 & -1 & -1 & 1 & 1 & 0 & 1 & 0 \\
7 & 5 & 0 & 0 & 0 & -1 & 1 & 0 & 1 & 0 \\
7 & 6 & -2 & -3 & -4 & 1 & 1 & 2 & 2 & 1 \\
7 & 7 & 0 & -1 & -1 & 0 & 0 & 0 & 2 & 0 \\
7 & 8 & 0 & 0 & -1 & 0 & 0 & 0 & 1 & 1 \\
\midrule
8 & 1 & 0 & -2 & -2 & 1 & 0 & 1 & 1 & 0 \\
8 & 2 & 0 & 0 & -1 & 0 & 0 & 0 & 1 & 1 \\
8 & 3 & 0 & 0 & 1 & 0 & 0 & 0 & -1 & 1 \\
8 & 4 & -2 & -3 & -4 & 2 & 1 & 1 & 1 & 2 \\
8 & 5 & 0 & -1 & -1 & 0 & 1 & 1 & 0 & 1 \\
8 & 6 & 0 & 0 & 0 & 0 & -1 & 1 & 0 & 1 \\
8 & 7 & 0 & -1 & 0 & 0 & 0 & 0 & 1 & 1 \\
8 & 8 & 0 & -1 & -1 & 0 & 0 & 0 & 0 & 2 \\
\bottomrule
\end{tabular}
\end{minipage}
\caption{Coefficients $c_k$ such that $L_i \cdot L_j = \sum_{k=1}^{8} c_k\, L_k$, for the bases defined above.}\label{tbl:appendix-LL}
\end{table}

\begin{table}[H]
\centering
{\bfseries Lemma~\ref{lem:L-sigmaLsbar}: $L \cdot \sigma(L\bar s) \subseteq \sigma(L\bar s)$.}\par\medskip
\footnotesize
\setlength{\tabcolsep}{4pt}
\begin{minipage}[t]{0.48\textwidth}
\centering
\begin{tabular}{cc|rrrrrrrr}
\toprule
$i$ & $j$ & $c_1$ & $c_2$ & $c_3$ & $c_4$ & $c_5$ & $c_6$ & $c_7$ & $c_8$ \\
\midrule
1 & 1 & -1 & -1 & -1 & 0 & 0 & 0 & 0 & 0 \\
1 & 2 & 0 & -2 & -2 & 0 & 1 & 1 & 0 & 1 \\
1 & 3 & 2 & 2 & 2 & 0 & -1 & -1 & 0 & -1 \\
1 & 4 & 0 & -1 & -2 & 0 & 1 & 0 & 1 & 0 \\
1 & 5 & 2 & 1 & 2 & -1 & -1 & 0 & 0 & -1 \\
1 & 6 & 0 & -2 & -2 & 1 & 0 & 0 & 1 & 1 \\
1 & 7 & 2 & 2 & 2 & -1 & -1 & -1 & -1 & 0 \\
1 & 8 & 2 & 1 & 2 & 0 & 0 & -1 & -1 & -1 \\
\midrule
2 & 1 & 2 & 2 & 2 & -1 & -1 & -1 & 0 & -1 \\
2 & 2 & -1 & 0 & -2 & 0 & 0 & 1 & 0 & 1 \\
2 & 3 & 1 & 1 & 2 & 0 & 0 & 0 & -1 & 0 \\
2 & 4 & 0 & 0 & -1 & 1 & 0 & 0 & 0 & 1 \\
2 & 5 & 0 & 0 & -1 & 0 & 1 & 1 & 0 & 0 \\
2 & 6 & 1 & 1 & 1 & 0 & -1 & 1 & -1 & 0 \\
2 & 7 & 1 & 2 & 2 & -1 & -1 & 0 & 0 & 0 \\
2 & 8 & 1 & 1 & 1 & -1 & 0 & 0 & -1 & 1 \\
\midrule
3 & 1 & 0 & -2 & -2 & 1 & 1 & 1 & 0 & 1 \\
3 & 2 & 1 & 2 & 2 & 0 & 0 & -1 & 0 & -1 \\
3 & 3 & -1 & -1 & 0 & 0 & 0 & 0 & 1 & 0 \\
3 & 4 & 0 & 0 & 1 & 1 & 0 & 0 & 0 & -1 \\
3 & 5 & 0 & 0 & 1 & 0 & 1 & -1 & 0 & 0 \\
3 & 6 & -1 & -1 & -1 & 0 & 1 & 1 & 1 & 0 \\
3 & 7 & -1 & -2 & -2 & 1 & 1 & 0 & 2 & 0 \\
3 & 8 & -1 & -1 & -1 & 1 & 0 & 0 & 1 & 1 \\
\midrule
4 & 1 & 2 & 1 & 2 & -1 & -1 & 0 & -1 & 0 \\
4 & 2 & 0 & 0 & -1 & 1 & 0 & 1 & 0 & 0 \\
4 & 3 & 0 & 0 & 1 & 1 & 0 & -1 & 0 & 0 \\
4 & 4 & 0 & -1 & -1 & 2 & 0 & 0 & 0 & 0 \\
4 & 5 & 2 & 2 & 3 & 0 & 0 & -1 & -1 & -1 \\
4 & 6 & 0 & -1 & 0 & 1 & 0 & 1 & 0 & 0 \\
4 & 7 & 0 & -1 & -1 & 1 & 0 & 0 & 1 & 1 \\
4 & 8 & 0 & 0 & 0 & 1 & 0 & 0 & -1 & 1 \\
\bottomrule
\end{tabular}
\end{minipage}\hfill
\begin{minipage}[t]{0.48\textwidth}
\centering
\begin{tabular}{cc|rrrrrrrr}
\toprule
$i$ & $j$ & $c_1$ & $c_2$ & $c_3$ & $c_4$ & $c_5$ & $c_6$ & $c_7$ & $c_8$ \\
\midrule
5 & 1 & 0 & -1 & -2 & 1 & 0 & 0 & 0 & 1 \\
5 & 2 & 0 & 1 & 0 & 0 & 1 & 0 & -1 & 0 \\
5 & 3 & 0 & -1 & 0 & 0 & 1 & 0 & 1 & 0 \\
5 & 4 & -2 & -3 & -4 & 2 & 2 & 1 & 1 & 1 \\
5 & 5 & 0 & -1 & -1 & 0 & 2 & 0 & 0 & 0 \\
5 & 6 & 0 & -1 & -1 & 0 & 1 & 1 & 0 & 1 \\
5 & 7 & 0 & 0 & -1 & 0 & 1 & 0 & 1 & 0 \\
5 & 8 & 0 & 0 & 0 & 0 & 1 & -1 & 0 & 1 \\
\midrule
6 & 1 & 2 & 1 & 2 & 0 & 0 & -1 & -1 & -1 \\
6 & 2 & 1 & 2 & 1 & -1 & 0 & 0 & 0 & -1 \\
6 & 3 & -1 & -1 & -1 & 1 & 0 & 1 & 1 & 0 \\
6 & 4 & 0 & 1 & 0 & 1 & 0 & 0 & 0 & -1 \\
6 & 5 & 0 & 0 & 0 & 0 & 1 & 0 & 1 & -1 \\
6 & 6 & -1 & -2 & -3 & 1 & 1 & 2 & 1 & 0 \\
6 & 7 & 1 & 1 & 1 & 0 & -1 & 0 & 1 & -1 \\
6 & 8 & 1 & 1 & 1 & 0 & 0 & 0 & 0 & 0 \\
\midrule
7 & 1 & 0 & -1 & -2 & 0 & 1 & 1 & 0 & 0 \\
7 & 2 & -1 & -1 & -2 & 0 & 1 & 1 & 1 & 1 \\
7 & 3 & 1 & 0 & 2 & 0 & -1 & 0 & 0 & 0 \\
7 & 4 & 0 & -1 & -1 & 1 & 0 & 1 & 1 & 0 \\
7 & 5 & 0 & -1 & 0 & 0 & 1 & 1 & 0 & 0 \\
7 & 6 & -1 & -2 & -2 & 0 & 0 & 2 & 1 & 1 \\
7 & 7 & 1 & 0 & 1 & -1 & 0 & 0 & 1 & 0 \\
7 & 8 & -1 & -3 & -3 & 1 & 1 & 1 & 1 & 2 \\
\midrule
8 & 1 & 0 & -2 & -2 & 1 & 0 & 1 & 1 & 0 \\
8 & 2 & -1 & -1 & -3 & 1 & 1 & 0 & 1 & 1 \\
8 & 3 & 1 & 2 & 3 & -1 & -1 & -1 & 0 & 0 \\
8 & 4 & 0 & 0 & 0 & 1 & 0 & -1 & 1 & 0 \\
8 & 5 & 0 & -1 & -1 & 0 & 1 & 0 & 1 & 1 \\
8 & 6 & -1 & -1 & -2 & 1 & 0 & 1 & 1 & 1 \\
8 & 7 & -1 & -1 & -2 & 0 & 0 & 0 & 2 & 1 \\
8 & 8 & 1 & 1 & 1 & 0 & -1 & -1 & 0 & 1 \\
\bottomrule
\end{tabular}
\end{minipage}
\caption{Coefficients $c_k$ such that $L_i \cdot N_j = \sum_{k=1}^{8} c_k\, N_k$, for the bases defined above.}\label{tbl:appendix-LN}
\end{table}

\begin{table}[H]
\centering
{\bfseries Lemma~\ref{lem:sigmaLs-closed}: $\sigma(Ls) \cdot \sigma(Ls) \subseteq \sigma(Ls)$.}\par\medskip
\footnotesize
\setlength{\tabcolsep}{4pt}
\begin{minipage}[t]{0.48\textwidth}
\centering
\begin{tabular}{cc|rrrrrrrr}
\toprule
$i$ & $j$ & $c_1$ & $c_2$ & $c_3$ & $c_4$ & $c_5$ & $c_6$ & $c_7$ & $c_8$ \\
\midrule
1 & 1 & -1 & 1 & 1 & 0 & 0 & 0 & 0 & 0 \\
1 & 2 & 0 & 0 & 2 & 0 & -1 & -1 & 0 & -1 \\
1 & 3 & -2 & -2 & -4 & 0 & 1 & 1 & 0 & 1 \\
1 & 4 & 0 & 1 & 2 & -2 & -1 & 0 & -1 & 0 \\
1 & 5 & -2 & -1 & -2 & 1 & -1 & 0 & 0 & 1 \\
1 & 6 & 0 & 2 & 2 & -1 & 0 & -2 & -1 & -1 \\
1 & 7 & -2 & -2 & -2 & 1 & 1 & 1 & -1 & 0 \\
1 & 8 & -2 & -1 & -2 & 0 & 0 & 1 & 1 & -1 \\
\midrule
2 & 1 & -2 & -3 & -2 & 0 & 1 & 1 & 0 & 1 \\
2 & 2 & 2 & -1 & 1 & 0 & 0 & 0 & 0 & 0 \\
2 & 3 & 2 & 3 & 3 & 0 & -2 & -2 & 0 & -2 \\
2 & 4 & 2 & 1 & 2 & -1 & 0 & -1 & 0 & -2 \\
2 & 5 & 4 & 4 & 7 & -2 & -3 & -2 & -1 & -2 \\
2 & 6 & 0 & -2 & -1 & 1 & 0 & -1 & 2 & 0 \\
2 & 7 & 2 & 1 & 2 & 0 & -1 & -2 & -1 & 0 \\
2 & 8 & 2 & 1 & 3 & 1 & -1 & -1 & -1 & -2 \\
\midrule
3 & 1 & 0 & 1 & 0 & 0 & -1 & -1 & 0 & -1 \\
3 & 2 & -2 & -3 & -5 & 0 & 2 & 2 & 0 & 2 \\
3 & 3 & 2 & 1 & 1 & 0 & 0 & 0 & 0 & 0 \\
3 & 4 & -2 & -3 & -6 & 1 & 2 & 1 & 2 & 2 \\
3 & 5 & 0 & -2 & -3 & 0 & 1 & 2 & 1 & 0 \\
3 & 6 & 0 & -2 & -3 & 1 & 0 & 1 & 0 & 2 \\
3 & 7 & 2 & 3 & 2 & -2 & -1 & 0 & -1 & 0 \\
3 & 8 & 2 & 1 & 1 & -1 & 1 & -1 & -1 & 0 \\
\midrule
4 & 1 & -2 & -1 & -2 & -1 & 1 & 0 & 1 & 0 \\
4 & 2 & 0 & -2 & -1 & -1 & 0 & 1 & 0 & 2 \\
4 & 3 & 4 & 4 & 5 & -1 & -2 & -1 & -2 & -2 \\
4 & 4 & 2 & 1 & 1 & -2 & 0 & 0 & 0 & 0 \\
4 & 5 & 2 & 2 & 3 & -2 & -2 & 1 & -1 & -1 \\
4 & 6 & 2 & 1 & 2 & -1 & -2 & -1 & 0 & 0 \\
4 & 7 & 4 & 5 & 7 & -3 & -2 & -2 & -3 & -1 \\
4 & 8 & 2 & 0 & 2 & -1 & 0 & 0 & -1 & -1 \\
\bottomrule
\end{tabular}
\end{minipage}\hfill
\begin{minipage}[t]{0.48\textwidth}
\centering
\begin{tabular}{cc|rrrrrrrr}
\toprule
$i$ & $j$ & $c_1$ & $c_2$ & $c_3$ & $c_4$ & $c_5$ & $c_6$ & $c_7$ & $c_8$ \\
\midrule
5 & 1 & 0 & 1 & 2 & -1 & -2 & 0 & 0 & -1 \\
5 & 2 & -2 & -5 & -6 & 2 & 1 & 2 & 1 & 2 \\
5 & 3 & 2 & 3 & 2 & 0 & -1 & -2 & -1 & 0 \\
5 & 4 & 0 & -1 & -2 & 0 & 0 & -1 & 1 & 1 \\
5 & 5 & 2 & 1 & 1 & 0 & -2 & 0 & 0 & 0 \\
5 & 6 & 0 & -1 & -1 & 2 & -1 & -1 & 0 & 1 \\
5 & 7 & 0 & 0 & -1 & 0 & -1 & 0 & -1 & 2 \\
5 & 8 & 2 & 2 & 2 & 0 & -1 & -1 & -2 & -1 \\
\midrule
6 & 1 & -2 & -2 & -2 & 1 & 0 & -1 & 1 & 1 \\
6 & 2 & 2 & 1 & 2 & -1 & 0 & -1 & -2 & 0 \\
6 & 3 & 2 & 3 & 2 & -1 & 0 & -1 & 0 & -2 \\
6 & 4 & 0 & 0 & -1 & -1 & 2 & -1 & 0 & 0 \\
6 & 5 & 2 & 2 & 2 & -2 & -1 & -1 & 0 & -1 \\
6 & 6 & 2 & 1 & 1 & 0 & 0 & -2 & 0 & 0 \\
6 & 7 & 2 & 4 & 3 & -1 & -1 & -2 & -2 & -1 \\
6 & 8 & 4 & 5 & 7 & -2 & -1 & -3 & -2 & -3 \\
\midrule
7 & 1 & 0 & 2 & 2 & -1 & -1 & -1 & -2 & 0 \\
7 & 2 & 0 & -2 & -1 & 0 & 1 & 2 & -1 & 0 \\
7 & 3 & 0 & -2 & -3 & 2 & 1 & 0 & 1 & 0 \\
7 & 4 & -2 & -4 & -6 & 1 & 2 & 2 & 1 & 1 \\
7 & 5 & 2 & 1 & 2 & 0 & -1 & 0 & -1 & -2 \\
7 & 6 & 0 & -3 & -2 & 1 & 1 & 0 & 0 & 1 \\
7 & 7 & 2 & 1 & 1 & 0 & 0 & 0 & -2 & 0 \\
7 & 8 & 0 & -1 & -1 & 1 & 2 & 0 & -1 & -1 \\
\midrule
8 & 1 & 0 & 1 & 2 & 0 & 0 & -1 & -1 & -2 \\
8 & 2 & 0 & -2 & -2 & -1 & 1 & 1 & 1 & 0 \\
8 & 3 & 0 & 0 & -2 & 1 & -1 & 1 & 1 & 0 \\
8 & 4 & 0 & 1 & -1 & -1 & 0 & 0 & 1 & -1 \\
8 & 5 & 0 & -1 & -1 & 0 & -1 & 1 & 2 & -1 \\
8 & 6 & -2 & -4 & -6 & 2 & 1 & 1 & 2 & 1 \\
8 & 7 & 2 & 2 & 2 & -1 & -2 & 0 & -1 & -1 \\
8 & 8 & 2 & 1 & 1 & 0 & 0 & 0 & 0 & -2 \\
\bottomrule
\end{tabular}
\end{minipage}
\caption{Coefficients $c_k$ such that $M_i \cdot M_j = \sum_{k=1}^{8} c_k\, M_k$, for the bases defined above.}\label{tbl:appendix-MM}
\end{table}

\subsection{The mod-2-quotient reading of Tables~\ref{tbl:appendix-LN}
and~\ref{tbl:appendix-MM}}\label{app:mod2quotient}

The integer tables above are the explicit certificates of
Lemmas~\ref{lem:L-sigmaLsbar} and~\ref{lem:sigmaLs-closed}; we now
record the structural account of \emph{what} they certify.

From $L\bar s \cap Ls = 2L$ (Section~\ref{subsec:wilson-leech}) one
has $2L \subseteq L\bar s$ and $2L \subseteq Ls$, and applying
$\sigma$ (which fixes $2L$),
\begin{equation}\label{eq:disp13}
  2L \subseteq \sigma(L\bar s), \qquad 2L \subseteq \sigma(Ls).
\end{equation}
Because $L \cdot L \subseteq L$ (Lemma~\ref{lem:L-order}) and the
octonion product is $\ZZ$-bilinear, $(a + 2c) \cdot b \equiv a \cdot
b \pmod{2L}$ for $c \in L$, and likewise in the right factor.  The
product therefore descends to a well-defined $\FF_2$-bilinear product
$\bar\mu$ on the \emph{mod-2 quotient}
\begin{equation}\label{eq:disp14}
  \bar L := L/2L \cong \FF_2^{\,8}, \qquad
  \bar\mu \colon \bar L \times \bar L \to \bar L,
\end{equation}
the reduction of the octonion order~$L$ modulo~$2$, an
eight-dimensional $\FF_2$-algebra.\footnote{It is not an octonion
algebra in the composition-algebra sense: those are simple, whereas
$\bar L$ has the proper two-sided ideal~$V$ exhibited below.}  The
structure constants of $\bar\mu$ are exactly the entries of
Table~\ref{tbl:appendix-LL} reduced modulo~$2$.

Since $2L$ lies in both $\sigma(L\bar s)$ and~$\sigma(Ls)$, those
sublattices are \emph{full preimages}, under $L \to \bar L$, of
$\FF_2$-subspaces
\begin{equation}\label{eq:disp15}
  V := \sigma(L\bar s)/2L, \qquad W := \sigma(Ls)/2L,
\end{equation}
and from $\sigma(L\bar s) + \sigma(Ls) = L$ and
$\sigma(L\bar s) \cap \sigma(Ls) = 2L$ one obtains $\bar L = V \oplus
W$ with $\dim V = \dim W = 4$.  Reducing the closure statements
modulo~$2L$ then gives
\begin{equation}\label{eq:disp16}
  \text{Lemma~\ref{lem:L-sigmaLsbar}} \iff
    \bar\mu(\bar L, V) \subseteq V, \qquad
  \text{Lemma~\ref{lem:sigmaLs-closed}} \iff
    \bar\mu(W, W) \subseteq W.
\end{equation}
That is: Table~\ref{tbl:appendix-LN} certifies that $V$ is a
\textbf{left ideal} of $\bar L$;
Table~\ref{tbl:appendix-MM} certifies that $W$ is a
\textbf{subalgebra} of it.  The integer tables and this mod-2-quotient
description are complementary: the tables verify, the quotient
explains.  The descent and both equivalences were checked in exact
arithmetic.\footnote{\scriptref{verify\_mod2\_quotient.py}.}

A note on terminology: the construction is a \emph{quotient}, the
surjection $L \to L/2L$, and is named the \emph{mod-2 quotient}.  It
is kept verbally distinct from a $2$-adic scaling or saturation, the
operations behind Corradetti's thematically adjacent construction
\cite{Corradetti2026}, which is contrasted with the present one in
Section~\ref{sec:related}.

\paragraph{Relation to prior work.}
The lattice-level picture of $\bar L \cong \FF_2^{\,8}$ as an
orthogonal space under the plus-type quadratic form induced by the
$E_8$ Gram form, together with its $4$-dimensional maximal totally
isotropic subspaces (Lagrangians), is Silverman's (part~12
of~\cite{Baez2014}); see Section~\ref{sec:related}.  The reading of
the Leech construction in terms of the preimages in~$L$ of a
\emph{complementary pair} of such Lagrangians is David Calderbank's,
in the discussion thread of part~9 \cite{Calderbank2023}.  Building
on those posts and their surrounding discussion, the present
construction equips $\bar L$ with the
$\FF_2$-bilinear product $\bar\mu$ inherited from the octonion
product $L \cdot L \subseteq L$, and identifies the two Lagrangians
of one specific complementary pair, $V = \sigma(L\bar s)/2L$ and
$W = \sigma(Ls)/2L$, as ideal and subalgebra of
$(\bar L, \bar\mu)$.

\paragraph{Polarization of $\bar L$ by $\sigma$.}
Direct verification yields two further structural facts about the
decomposition $\bar L = V \oplus W$.\footnote{\scriptref{verify\_discovery1\_W\_subalgebra.py}
\scriptref{verify\_discovery2\_V\_isotropic.py}.}  First, $V$ is in fact a
\emph{two-sided} ideal of $\bar L$: the right-ideal direction
$\bar\mu(V, \bar L) \subseteq V$ holds in addition to the left-ideal
direction (Lemma~\ref{lem:L-sigmaLsbar}).  The right-ideal direction
is not formally implied by the left-ideal direction, since $\bar\mu$
is non-commutative.\footnote{Commutativity is a property of the
algebra, not of a basis; what matters is \emph{which order} is being
reduced.  Reducing the Gravesian order
$\ZZ\langle e_0,\ldots,e_7\rangle$ gives a commutative
$\FF_2$-algebra, since $e_i \cdot e_j - e_j \cdot e_i =
2\,e_i \cdot e_j$ lies in twice that order.  The order~$L$
also contains half-integer elements such as~$s$, and for these the
commutator $a \cdot b - b \cdot a$ need not lie in~$2L$; direct
computation confirms that $\bar\mu$ is non-commutative.}

Second, equip $\bar L$ with the natural plus-type quadratic form
\begin{equation}\label{eq:disp17}
  q \colon \bar L \to \FF_2, \qquad q(v + 2L) \;:=\; \Nrm(v)/2 \pmod 2,
\end{equation}
well-defined because $L$ is even
($\Nrm(v + 2u) \equiv \Nrm(v) \pmod 4$).  For $L = D_8^+$ this is the
standard plus-type form on~$\FF_2^{\,8}$ with Witt index~$4$.  Direct
verification shows that \emph{both}~$V$ and~$W$ are totally isotropic
under~$q$ (every element of either subspace has $q = 0$); each is
therefore a maximal totally isotropic subspace, a
\emph{Lagrangian}.  The decomposition $\bar L = V \oplus W$ is thus a
\emph{Witt decomposition of $(\bar L, q)$ into a complementary pair
of Lagrangians}, a polarization of the orthogonal $\FF_2$-space.

\paragraph{The structural reframing.}
In this picture, Lemmas~\ref{lem:L-sigmaLsbar}
and~\ref{lem:sigmaLs-closed} together express a single fact:
\emph{the construction of Section~\ref{sec:construction} produces a
polarization of $\bar L$}, with one Lagrangian carrying the
two-sided-ideal role ($V$) and the other carrying the subalgebra
role ($W$).  The open question recorded in
Section~\ref{sec:conclusion}, a structural reason for
Lemma~\ref{lem:sigmaLs-closed}, accordingly opens into a research
program: which polarizations of $\bar L$ arise from this
construction, and whether any such polarization lifts to a closed
bilinear product on $\LL$ itself.

% ===============================================================
\section{Historical note: integral Cayley numbers,
1923--1946}\label{app:history}
% ===============================================================

The construction this paper builds on was developed across four
primary papers between 1923 and 1946.  Each was obtained in the
original and re-derived in exact arithmetic in this
project.\footnote{\scriptref{verify\_dickson\_1923.py}
\scriptref{verify\_kirmse\_1924.py}
\scriptref{verify\_kirmse\_1924\_forensic.py} \scriptref{verify\_mahler\_1942.py}
\scriptref{verify\_coxeter\_1946.py} Verification notes accompany each,
in \texttt{paper/reviews/}.}

\noindent\textbf{Dickson (1923) \cite{Dickson1923}} gave the first
construction: the Cayley--Dickson doubled product and an explicit
8-element basis of an order $O_D$, with all 64 basis products
closing.  His Theorem~XV states three maximal orders; the true count
is \emph{seven}.  The cause is in his derivation (p.~321): he assumes
the orders contain the Hurwitz quaternions, and exactly three of the
seven do.

\noindent\textbf{Kirmse (1924) \cite{Kirmse1924}} took up the problem
independently.  His multiplication table is itself a correct
octonion algebra, and he derived a correct \emph{necessary} form
for the maximal orders, without proving sufficiency; thirty
$E_8$-type lattices realize the form.  Among the eight of those lattices that
are stable under multiplication by units (an unstated narrowing of
the criterion, in our reconstruction), he exhibited the one,
$J_1$ (p.~70), that is not an order.  The remaining seven of his
eight are the genuine maximal orders.  Kirmse also develops an
ideal-theoretic treatment of the integral octonions, a substantive
companion contribution that is taken up again in Mahler (1942),
and methodically takes algebra alternativity as a foundational
constraint, anticipating the structural role alternativity plays in
later integral-octonion arithmetic.

\noindent\textbf{Mahler (1942) \cite{Mahler1942}} developed the
arithmetic of a maximal order: for every Cayley number~$X$ there is
an integral~$G$ with $\Nrm(X-G)\le\tfrac12$ (sharper than Dickson's
$5/4$), and every left or right ideal is principal, generated by an
integral Cayley number of norm $1$, $2$, or $4$.

\noindent\textbf{Coxeter (1946) \cite{Coxeter1946}} identifies
Kirmse's error with the same non-closure witness and records that
the remedy is due to Bruck: ``Bruck's domain $J$ can be
derived from $J_1$ by transposing two of the $i$'s'' (\S4).  This is
a transposition of two imaginary basis units, a linear involution
of~$\RR^8$; the $\binom72 = 21$ such transpositions yield the seven
maximal orders in seven sets of three.  His \S\S6--13 develop the
$E_8/5_{21}$ geometry of the corrected order.  A postscript (\S14)
records his after-the-fact awareness of Dickson and Mahler and
identifies Dickson's undercount with exactly the cause above.

It is curious to find that the Coxeter--Bruck transposition would
come in useful here: the same linear involution of~$\RR^8$, acting
now on a different object.  The Coxeter--Bruck transposition acts
on a maximal-order candidate that sits asymmetrically in~$\RR^8$;
the present~$\sigma$ acts on the coordinate-symmetric placement
$L = D_8^+$ (where it is inert at the $E_8$ stratum) and does its
work on the sublattice~$Ls$ of Wilson's Leech construction.

% ===============================================================
\section{Research methodology}\label{sec:methodology}
% ===============================================================

This research was conducted as a collaboration between a human
researcher and an AI agent (Claude, Anthropic) over a series of
directed prompts, governed by a written methodology (the
``Manifesto'') established at the outset.

The human researcher (the author) brought the initial hypothesis
(motivated by prior work on Okubo algebras and autotopies of normed
composition algebras), established the research methodology, and
maintained control of the inquiry through prompt-by-prompt review.
The AI agent performed the computational work: implementing octonion
algebras, Wilson's Leech lattice construction, trial algebras, and
closure tests, and contributed mathematical analysis.

Tests were written from first principles to independently verify
every foundation: that octonion multiplication satisfies the
composition property $\Nrm(xy) = \Nrm(x)\Nrm(y)$, that Wilson's
construction produces a lattice with the correct minimal shell of
196{,}560 vectors, and that the $E_8$ lattice is indeed a maximal
order.  This independent verification is how we know that the
properties stated in the literature are correctly implemented.
Closure of the twisted product was first observed on over
12{,}000{,}000 random pairs of minimal vectors with zero failures,
which prompted the search for a symbolic proof.

\paragraph{Symbolic proof and numeric validation.}
The proof of Theorem~\ref{thm:main} given in Section~\ref{sec:proof}
is symbolic: three of the four lemmas are established by exhibiting a
finite list of integer coefficients
(Appendix~\ref{sec:appendix-tables}); the fourth, $\sigma(L) = L$,
is a parity argument.  The conclusion follows by $\ZZ$-bilinearity.  No floating-point
or sampled arithmetic enters the proof.  Independently, the same
basis-by-basis identities that constitute the appendix tables were
recomputed in two separate computer-algebra implementations as a
cross-check: a Python implementation using
\texttt{fractions.Fraction} exact rationals, and a GAP/LOOPS
implementation using GAP's native exact rational
arithmetic.\footnote{\scriptref{symbolic\_proof\_checks.py}
\gapref{tests/test\_lemmas.g}.}  Both implementations
verify the same basis-pair identities as the appendix tables.
This redundancy is validation, not proof: it serves to detect
transcription or implementation errors in the appendix tables, but
the logical content of the proof rests on the tables themselves.

The methodology required:
\begin{itemize}
\item Every claim backed by a reference or computation.
\item Every prompt logged, with actions recorded.
\item Every conclusion independently verified.
\item Forward-evolving corrections: no rewriting history; errors are
  corrected in subsequent entries.
\end{itemize}

The path to the result was not direct.  Six trial algebra families
were explored and ruled out before the transposition twist was
identified in trial~007.  The systematic methodology was essential:
understanding the anatomy of failure (which conditions fail, for
which vector types, and why) before trying the next approach.  The failure analysis of the untwisted product (Wilson's
condition~(3) per Definition~\ref{def:wilson-leech}, exclusively
type-3$\times$type-3) directly pointed to the
mechanism that the twist needed to correct.

Iterative convergence toward a robust result, including learning
along the way what it is we are actually doing, is how study,
research, and discovery have always proceeded.  Mistakes were made
on both sides: the AI agent occasionally departed from the agreed
protocol or produced incorrect intermediate analyses, and the human
author at times framed prompts under misconceptions or accepted
intermediate results without sufficient scrutiny.  In each case, the
error was caught by the systematic cross-checking that the
methodology requires (explicit verification, prompt-by-prompt
review, and independent reimplementation in two computer-algebra
systems) and corrected forward.  What is novel here is the
ability to trace the entire path, beginning to end, through the
complete list of prompts.

The full record is available at the repository~\cite{repo}: all
prompts, code, evidence, and git history.  The reader can form
their own assessment of what this means for questions of discovery
and attribution.

% ===============================================================
% References
% ===============================================================

\begin{thebibliography}{XXXXXXXX}

\bibitem[Baez2002]{Baez2002}
J.\,C.~Baez,
\emph{The octonions},
Bull.\ Amer.\ Math.\ Soc.\ \textbf{39} (2002), 145--205.
\href{https://arxiv.org/abs/math/0105155}{arXiv:math/0105155}.

\bibitem[Baez2014]{Baez2014}
J.\,C.~Baez (with contributions by G.~Egan, T.~Silverman and others),
\emph{Integral octonions} (parts 1--12),
The $n$-Category Caf\'e, 2013--2016.
The material discussed here appears in parts 9--12: Egan's Leech
construction (part~9, 23 November 2014) and the Jordan-ring count
(part~10, 1 December 2014),
\url{https://golem.ph.utexas.edu/category/2014/11/integral_octonions_part_9.html},
\url{https://golem.ph.utexas.edu/category/2014/12/integral_octonions_part_10.html};
the 270 rescaled copies of the first $E_8$ shell (part~11, 12 December
2014) and Tim Silverman's account of the integral octonions modulo~2
(part~12, 31 January 2016),
\url{http://math.ucr.edu/home/baez/octonions/integers/integers_11.html},
\url{http://math.ucr.edu/home/baez/octonions/integers/integers_12.html}.

\bibitem[Calderbank2023]{Calderbank2023}
D.~Calderbank,
comment of 2 April 2023 on Baez~\cite{Baez2014}, part~9,
\url{https://golem.ph.utexas.edu/category/2014/11/integral_octonions_part_9.html}.

\bibitem[ConwaySloane1982]{ConwaySloane1982}
J.\,H.~Conway and N.\,J.\,A.~Sloane,
\emph{Twenty-three constructions for the Leech lattice},
Proc.\ Roy.\ Soc.\ London Ser.\ A \textbf{381} (1982), 275--283.

\bibitem[ConwaySloane1999]{ConwaySloane1999}
J.\,H.~Conway and N.\,J.\,A.~Sloane,
\emph{Sphere Packings, Lattices and Groups},
3rd edition, Springer-Verlag, 1999.

\bibitem[ConwaySmith2003]{ConwaySmith2003}
J.\,H.~Conway and D.\,A.~Smith,
\emph{On Quaternions and Octonions: Their Geometry, Arithmetic, and
Symmetry},
A.\,K.~Peters, 2003.

\bibitem[Corradetti2026]{Corradetti2026}
D.~Corradetti,
\emph{Integral elements of Okubo algebra and the $E_8$-lattice},
preprint, 2026.
\href{https://arxiv.org/abs/2605.09333}{arXiv:2605.09333}.

\bibitem[Coxeter1946]{Coxeter1946}
H.\,S.\,M.~Coxeter,
\emph{Integral Cayley numbers},
Duke Math.\ J.\ \textbf{13} (1946), 561--578.

\bibitem[Dickson1923]{Dickson1923}
L.\,E.~Dickson,
\emph{A new simple theory of hypercomplex integers},
J.~Math.\ Pures Appl.\ (9) \textbf{2} (1923), 281--326.

\bibitem[Dixon1995]{Dixon1995}
G.\,M.~Dixon,
\emph{Octonions: invariant Leech lattice exposed},
preprint, 1995.
\href{https://arxiv.org/abs/hep-th/9506080}{arXiv:hep-th/9506080}.

\bibitem[Dixon2005]{Dixon2005}
G.\,M.~Dixon,
\emph{Integral octonions, octonion multiplication and the Leech
lattice}, preprint, 10 March 2005.
\href{http://www.7stones.com/Homepage/8Lattice.pdf}{www.7stones.com/Homepage/8Lattice.pdf};
reference~[12] of~\cite{Wilson2009}.

\bibitem[Dixon2010]{Dixon2010}
G.\,M.~Dixon,
\emph{Integral octonions, octonion XY-product, and the Leech lattice},
preprint, 2010.
\href{https://arxiv.org/abs/1011.2541}{arXiv:1011.2541}.

\bibitem[Elduque1996]{Elduque1996}
A.~Elduque,
\emph{On a class of ternary composition algebras},
J.~Korean Math.\ Soc.\ \textbf{33} (1996), no.~1, 183--203.

\bibitem[Elduque2000\_Triality]{Elduque2000_Triality}
A.~Elduque,
\emph{On triality and automorphisms and derivations of composition
algebras},
Linear Algebra Appl.\ \textbf{314} (2000), no.~1--3, 49--74.

\bibitem[Elduque2023\_IsotropicNorm]{Elduque2023_IsotropicNorm}
A.~Elduque,
\emph{Okubo algebras with isotropic norm},
in: Non-Associative Algebras and Related Topics (NAART II, 2022),
Springer Proc.\ Math.\ Stat., vol.~427, Springer, Cham, 2023,
pp.~287--301.
\href{https://arxiv.org/abs/2210.03456}{arXiv:2210.03456}.

\bibitem[ElkiesGross1996]{ElkiesGross1996}
N.\,D.~Elkies and B.\,H.~Gross,
\emph{The exceptional cone and the Leech lattice},
Internat.\ Math.\ Res.\ Notices \textbf{1996}, no.~14, 665--698.

\bibitem[FureyHughes2025\_TrioOfTrialities]{FureyHughes2025_TrioOfTrialities}
N.~Furey and M.\,J.~Hughes,
\emph{Three generations and a trio of trialities},
Phys.\ Lett.\ B \textbf{865} (2025), 139473.
\href{https://arxiv.org/abs/2409.17948}{arXiv:2409.17948}.

\bibitem[GresnigtGourlayVarma2023]{GresnigtGourlayVarma2023}
N.\,G.~Gresnigt, L.~Gourlay, and A.~Varma,
\emph{Three generations of colored fermions with $S_3$ family
symmetry from Cayley--Dickson sedenions},
Eur.\ Phys.\ J.\ C \textbf{83} (2023), 747.
\href{https://arxiv.org/abs/2306.13098}{arXiv:2306.13098}.

\bibitem[Hall2019\_Moufang]{Hall2019_Moufang}
J.\,I.~Hall,
\emph{Moufang Loops and Groups with Triality are Essentially the
Same Thing},
Memoirs of the American Mathematical Society, vol.~260, no.~1252,
Amer.\ Math.\ Soc., Providence, RI, 2019.

\bibitem[KamiyaOkubo2015]{KamiyaOkubo2015}
N.~Kamiya and S.~Okubo,
\emph{Algebras satisfying triality and\/ $S_4$-symmetry},
preprint, 2015.
\href{https://arxiv.org/abs/1503.00614}{arXiv:1503.00614}.

\bibitem[Kirmse1924]{Kirmse1924}
J.~Kirmse,
\emph{\"Uber die Darstellbarkeit nat\"urlicher ganzer Zahlen als Summen
von acht Quadraten und \"uber ein mit diesem Problem zusammenh\"angendes
nichtkommutatives und nichtassoziatives Zahlensystem},
Ber.\ Verh.\ S\"achs.\ Akad.\ Wiss.\ Leipzig, Math.-Phys.\ Kl.\
\textbf{76} (1924), 63--82.

\bibitem[Koeplinger2023]{Koeplinger2023}
J.~K\"oplinger,
\emph{Towards autotopies of normed composition algebras in algebraic
quantum field theory},
preprint, 2023.
\href{https://doi.org/10.13140/RG.2.2.19003.39209}{DOI:\,10.13140/RG.2.2.19003.39209}.

\bibitem[LepowskyMeurman1982]{LepowskyMeurman1982}
J.~Lepowsky and A.~Meurman,
\emph{An\/ $E_8$-approach to the Leech lattice and the Conway group},
J.~Algebra \textbf{77} (1982), no.~2, 484--504.

\bibitem[Mahler1942]{Mahler1942}
K.~Mahler,
\emph{On ideals in the Cayley--Dickson algebra},
Proc.\ Roy.\ Irish Acad.\ Sect.\ A \textbf{48} (1942), 123--133.

\bibitem[MarraniCorradettiZucconi2025]{MarraniCorradettiZucconi2025}
A.~Marrani, D.~Corradetti, and F.~Zucconi,
\emph{Physics with non-unital algebras? An invitation to the Okubo
algebra},
J.\ Phys.\ A: Math.\ Theor.\ \textbf{58} (2025), 075202.
\href{https://arxiv.org/abs/2309.17435}{arXiv:2309.17435}.

\bibitem[Okubo1995\_Book]{Okubo1995_Book}
S.~Okubo,
\emph{Introduction to Octonion and Other Non-Associative Algebras in
Physics},
Montroll Memorial Lecture Series in Mathematical Physics, vol.~2,
Cambridge University Press, Cambridge, 1995.

\bibitem[Petersson2018]{Petersson2018}
H.\,P.~Petersson,
\emph{Integral octonions},
Lecture, M\'alaga Workshop on Non-Associative Algebras, 2018.
\url{https://www.fernuni-hagen.de/mi/fakultaet/emeriti/docs/petersson/ass.-rem.-int.-oct.pdf}.

\bibitem[SmithVojtechovsky2022]{SmithVojtechovsky2022}
J.\,D.\,H.~Smith and P.~Vojt\v{e}chovsk\'y,
\emph{Okubo quasigroups},
Doc.\ Math.\ \textbf{27} (2022), 535--580.

\bibitem[Wilson2009]{Wilson2009}
R.\,A.~Wilson,
\emph{Octonions and the Leech lattice},
J.~Algebra \textbf{322} (2009), 2186--2190.

\bibitem[repo]{repo}
J.~K\"oplinger,
\emph{An order on the Leech lattice from a $\ZZ_3$-symmetric
triple-octonion product: research repository},
\url{https://bitbucket.org/jenskoeplinger/leech_alg} and
\url{https://github.com/koeplinger/leech_alg}, 2026.

\end{thebibliography}

\end{document}
